% 康奈尔听课笔记模板（配色取自王道课件）
\documentclass[10.5pt,a4paper,oneside]{article}
\usepackage[a4paper,top=2.3cm,bottom=2.2cm,left=5.7cm,right=1.5cm,
  marginparwidth=4.25cm,marginparsep=0.55cm,headheight=14pt,footskip=1.1cm]{geometry}
\usepackage{xeCJK}
\usepackage{xeCJKfntef}
\xeCJKDeclareCharClass{FullLeft}{"201C,"2018}
\xeCJKDeclareCharClass{FullRight}{"201D,"2019}
\xeCJKDeclareCharClass{CJK}{"2460->"24FF,"2605,"2606}
\setmainfont{Noto Sans CJK SC}[BoldFont={Noto Sans CJK SC Bold},ItalicFont={Noto Sans CJK SC},ItalicFeatures={FakeSlant=0.18}]
\setsansfont{Noto Sans CJK SC}[BoldFont={Noto Sans CJK SC Bold}]
\setmonofont{DejaVu Sans Mono}[Scale=0.9]
\setCJKmainfont{Noto Sans CJK SC}[BoldFont={Noto Sans CJK SC Bold},ItalicFont={Noto Sans CJK SC},ItalicFeatures={FakeSlant=0.18}]
\setCJKsansfont{Noto Sans CJK SC}[BoldFont={Noto Sans CJK SC Bold}]
\setCJKmonofont{Noto Sans Mono CJK SC}[Scale=0.9]
\linespread{1.32}
\usepackage{amsmath,amssymb}
\usepackage[dvipsnames,table]{xcolor}
\usepackage{graphicx}
\usepackage{adjustbox}
\usepackage{tikz}
\usetikzlibrary{calc}
\usepackage[most]{tcolorbox}
\usepackage{eso-pic}
\usepackage{fancyhdr}
\usepackage{titlesec}
\usepackage{needspace}
\usepackage{enumitem}
\usepackage{xltabular,booktabs,array,colortbl}
\usepackage{marginnote}
\reversemarginpar
\renewcommand*{\raggedleftmarginnote}{}
\renewcommand*{\raggedrightmarginnote}{}
\usepackage{caption}
\usepackage{titletoc}
\usepackage{fancyvrb}
\usepackage{framed}
\usepackage[section]{placeins}
\usepackage{hyperref}
\hypersetup{colorlinks=true,linkcolor=wdblue,urlcolor=wdblue,pdftitle={3.3
课后习题：栈和队列的应用},pdfauthor={听课笔记}}

% ---------- 课件配色 ----------
\definecolor{wdblue}{HTML}{4472C4}
\definecolor{wdlightblue}{HTML}{5B9BD5}
\definecolor{wdtitle}{HTML}{4A86C8}
\definecolor{wddark}{HTML}{1F3864}
\definecolor{wdorange}{HTML}{ED7D31}
\definecolor{wdgray}{HTML}{A5A5A5}
\definecolor{wdgold}{HTML}{FFC000}
\definecolor{wdgreen}{HTML}{70AD47}
\definecolor{wddeepgreen}{HTML}{3F7F2A}
\definecolor{wdred}{HTML}{C00000}
\definecolor{wdpurple}{HTML}{7030A0}
\definecolor{cuebg}{HTML}{F2F6FC}
\definecolor{codebg}{HTML}{F7F9FC}
\color[HTML]{262626}

% ---------- 行内样式 ----------
\newcommand{\hlred}[1]{\textcolor{wdred}{\bfseries #1}}
\newcommand{\hlblue}[1]{\textcolor{wdblue}{\bfseries #1}}
\newcommand{\hlgreen}[1]{\textcolor{wddeepgreen}{\bfseries #1}}
\newcommand{\hlorange}[1]{\textcolor{wdorange}{\bfseries #1}}
\newcommand{\hlpurple}[1]{\textcolor{wdpurple}{\bfseries #1}}
\newcommand{\uln}[1]{\CJKunderline[format=\color{wdred},thickness=0.9pt,depth=0.28em]{#1}}
\newcommand{\wavy}[1]{\CJKunderwave[format=\color{wdorange}]{#1}}
\newcommand{\hlmark}[1]{\CJKunderline[format=\color{wdgold!55},thickness=0.55em,depth=-0.05em]{#1}}
\newcommand{\tagbox}[2]{\tcbox[on line,colback=#1,colframe=#1,coltext=white,boxsep=0.5pt,left=3pt,right=3pt,top=0.5pt,bottom=0.5pt,arc=2pt,fontupper=\small\bfseries]{#2}}

% ---------- 康奈尔版式背景：左侧线索栏 ----------
\newlength{\cuestripw}\setlength{\cuestripw}{5.15cm}
\AddToShipoutPictureBG{%
  \begin{tikzpicture}[remember picture,overlay]
    \fill[cuebg] (current page.north west) rectangle ($(current page.south west)+(\cuestripw,0)$);
    \draw[wdblue,line width=1.2pt] ($(current page.north west)+(\cuestripw,-1.55cm)$) -- ($(current page.south west)+(\cuestripw,1.55cm)$);
    \node[anchor=west,text=wdblue,font=\small\bfseries] at ($(current page.north west)+(0.75cm,-1.25cm)$) {线索 · 问题};
    \node[anchor=west,text=wdblue,font=\small\bfseries] at ($(current page.north west)+(\cuestripw+0.55cm,-1.25cm)$) {笔记 · 老师主线 + 补充讲解};
  \end{tikzpicture}}

% ---------- 页眉页脚 ----------
\pagestyle{fancy}\fancyhf{}
\renewcommand{\headrulewidth}{0pt}
\fancyhead[R]{\small\color{wdgray}第3章 栈、队列和数组}
\fancyfoot[C]{\small\color{wdgray}\thepage}
\fancyfoot[R]{\small\color{wdorange}3.3
课后习题：栈和队列的应用 · 听课笔记}
\fancyfoot[L]{\small\hyperlink{toc}{\textcolor{wdblue}{\ensuremath{\uparrow}\,回到目录}}}
\fancypagestyle{plain}{\fancyhf{}\fancyfoot[C]{\small\color{wdgray}\thepage}\fancyfoot[R]{\small\color{wdorange}3.3
课后习题：栈和队列的应用 · 听课笔记}}

% ---------- 标题：仿课件的蓝色标题 + 六色条 ----------
\newcommand{\sixbar}{\par\nointerlineskip\vspace{3pt}\centering
  \begin{tikzpicture}\foreach \c [count=\i] in {wdlightblue,wdorange,wdgray,wdgold,wdblue,wdgreen}
     {\fill[\c] ({(\i-1)*1.05},0) rectangle ({(\i-1)*1.05+0.92},0.11);}\end{tikzpicture}\par}
\titleformat{\section}[block]{\centering\Large\bfseries\color{wdtitle}}{}{0pt}{}[\sixbar]
\titlespacing*{\section}{0pt}{16pt}{10pt}
\pretocmd{\section}{\Needspace{8\baselineskip}}{}{}
\titleformat{\subsection}{\Needspace{7\baselineskip}\large\bfseries\color{wddark}}{\textcolor{wdblue}{\rule[-0.15em]{4pt}{1.05em}}\hspace{6pt}}{0pt}{}
\titlespacing*{\subsection}{0pt}{12pt}{5pt}
\titleformat{\subsubsection}{\Needspace{5\baselineskip}\normalsize\bfseries\color{wdblue}}{\textcolor{wdorange}{\large\ensuremath{\blacktriangleright}}\hspace{4pt}}{0pt}{}
\titlespacing*{\subsubsection}{0pt}{8pt}{3pt}
\setcounter{secnumdepth}{0}
\setlength{\parindent}{0pt}\setlength{\parskip}{5pt}
\setlist{nosep,leftmargin=1.5em,itemsep=2pt}
\setlist[itemize,1]{label=\textcolor{wdblue}{\small\ensuremath{\bullet}}}
\setlist[itemize,2]{label=\textcolor{wdorange}{\scriptsize\ensuremath{\blacktriangleright}}}
\setlist[enumerate,1]{label=\textcolor{wdblue}{\bfseries\arabic*.}}
\providecommand{\tightlist}{\setlength{\itemsep}{2pt}\setlength{\parskip}{0pt}}

% ---------- 目录 ----------
\setcounter{tocdepth}{1}
\titlecontents{section}[0pt]{\addvspace{5pt}\bfseries\color{wddark}}{}{}{\titlerule*[0.7pc]{\color{wdgray!60}.}\contentspage}
\titlecontents{subsection}[1.6em]{\addvspace{1pt}\small\color{black!75}}{}{}{\titlerule*[0.7pc]{\color{wdgray!40}.}\contentspage}
\newcommand{\tocback}{\par\vspace{-3pt}{\hfill\footnotesize\hyperlink{toc}{\textcolor{wdgray}{\ensuremath{\uparrow}\,回到目录}}}\par\vspace{-4pt}}

% ---------- 线索栏 ----------
\newcommand{\cueblock}[2]{\Needspace{4\baselineskip}\leavevmode#1{#2}\par\vspace{-\baselineskip}\vspace{-\parskip}}
% 线索直接贴在段首行的左侧（不用 marginnote，避免分页时偶尔落进正文栏）
\newcommand{\cuebox}[1]{\leavevmode\llap{\smash{\begin{minipage}[t]{\marginparwidth}\raggedright\small\linespread{1.18}\selectfont #1\end{minipage}}\hspace{\marginparsep}}}
\newcommand{\cuenote}[1]{\cuebox{{\color{wdblue}\bfseries Q\,}\color{wddark}#1}}
\newcommand{\cuekey}[1]{\cuebox{\color{wdred}\bfseries #1}}

% ---------- 讲解框 ----------
\tcbset{wdbox/.style={enhanced,breakable,arc=2pt,boxrule=0pt,leftrule=3pt,left=7pt,right=7pt,top=5pt,bottom=5pt,
   before skip=7pt,after skip=7pt,fonttitle=\bfseries\small,attach title to upper={\quad},coltitle=#1,
   colframe=#1,colback=#1!6}}
\newtcolorbox{defbox}[1][]{enhanced,breakable,arc=3pt,boxrule=0.8pt,colframe=wdblue,colback=wdblue!4,
   left=8pt,right=8pt,top=6pt,bottom=6pt,fonttitle=\bfseries\small,coltitle=white,colbacktitle=wdblue,
   title={定义#1},before skip=8pt,after skip=8pt}
\newtcolorbox{thmbox}[1]{enhanced,breakable,arc=3pt,boxrule=0.8pt,colframe=wdblue,colback=wdblue!4,
   left=8pt,right=8pt,top=6pt,bottom=6pt,fonttitle=\bfseries\small,coltitle=white,colbacktitle=wdblue,
   title={#1},before skip=8pt,after skip=8pt}
\newtcolorbox{bridgebox}{wdbox=wdblue,title={承上启下}}
\newtcolorbox{plainbox}{wdbox=wddeepgreen,colframe=wdgreen,colback=wdgreen!8,title={大白话}}
\newtcolorbox{deepbox}{wdbox=wdpurple,title={背后的思想}}
\newtcolorbox{exambox}{wdbox=wdred,title={考点}}
\newtcolorbox{pitfallbox}{wdbox=wdorange,title={易错提醒}}
\newtcolorbox{bookbox}{wdbox=wddark,colback=wdgray!10,colframe=wdgray,title={王道书补充}}
\newtcolorbox{teacherbox}{wdbox=wdorange,colback=wdgold!8,colframe=wdgold,coltitle=wdorange,title={老师原话}}
\newtcolorbox{fixbox}{wdbox=wdred,colback=wdred!4,title={勘误 / 说明}}
% 习题笔记用
\newtcolorbox{questionbox}[1][]{enhanced,breakable,arc=3pt,boxrule=0.8pt,colframe=wddark,colback=white,
   left=8pt,right=8pt,top=6pt,bottom=6pt,fonttitle=\bfseries\small,coltitle=white,colbacktitle=wddark,
   title={#1},before skip=10pt,after skip=6pt}
\newtcolorbox{videobox}{wdbox=wdorange,colback=wdgold!8,colframe=wdgold,coltitle=wdorange,title={视频讲解}}
\newtcolorbox{jiexibox}{wdbox=wddark,colback=wdgray!10,colframe=wdgray,title={书上解析}}
\newtcolorbox{addbox}{wdbox=wdpurple,title={我的补充}}
\newtcolorbox{recallbox}{wdbox=wdblue,title={回顾正课}}
\newtcolorbox{watchbox}{enhanced,breakable,arc=3pt,boxrule=1pt,colframe=wdorange,colback=wdgold!12,
   fonttitle=\bfseries\small,coltitle=white,colbacktitle=wdorange,title={建议去看视频},left=8pt,right=8pt,before skip=8pt,after skip=8pt}
\newtcolorbox{summarybox}{enhanced,arc=3pt,
   boxrule=1pt,colframe=wdblue,colback=white,colbacktitle=wdblue,coltitle=white,fonttitle=\bfseries,
   title={总结区 · 用自己的话复述这一节},left=10pt,right=10pt,top=8pt,bottom=8pt,before skip=14pt}

% ---------- 代码 ----------
\usepackage{color}
\usepackage{fancyvrb}
\newcommand{\VerbBar}{|}
\newcommand{\VERB}{\Verb[commandchars=\\\{\}]}
\DefineVerbatimEnvironment{Highlighting}{Verbatim}{commandchars=\\\{\}}
% Add ',fontsize=\small' for more characters per line
\newenvironment{Shaded}{}{}
\newcommand{\AlertTok}[1]{\textcolor[rgb]{1.00,0.00,0.00}{#1}}
\newcommand{\AnnotationTok}[1]{\textcolor[rgb]{0.38,0.63,0.69}{#1}}
\newcommand{\AttributeTok}[1]{\textcolor[rgb]{0.49,0.56,0.16}{#1}}
\newcommand{\BaseNTok}[1]{\textcolor[rgb]{0.77,0.35,0.07}{#1}}
\newcommand{\BuiltInTok}[1]{\textcolor[rgb]{0.00,0.50,0.00}{#1}}
\newcommand{\CharTok}[1]{\textcolor[rgb]{0.63,0.19,0.43}{#1}}
\newcommand{\CommentTok}[1]{\textcolor[rgb]{0.18,0.55,0.23}{#1}}
\newcommand{\CommentVarTok}[1]{\textcolor[rgb]{0.18,0.55,0.23}{#1}}
\newcommand{\ConstantTok}[1]{\textcolor[rgb]{0.77,0.35,0.07}{#1}}
\newcommand{\ControlFlowTok}[1]{\textcolor[rgb]{0.12,0.31,0.71}{#1}}
\newcommand{\DataTypeTok}[1]{\textcolor[rgb]{0.12,0.31,0.71}{#1}}
\newcommand{\DecValTok}[1]{\textcolor[rgb]{0.77,0.35,0.07}{#1}}
\newcommand{\DocumentationTok}[1]{\textcolor[rgb]{0.18,0.55,0.23}{#1}}
\newcommand{\ErrorTok}[1]{\textcolor[rgb]{1.00,0.00,0.00}{#1}}
\newcommand{\ExtensionTok}[1]{#1}
\newcommand{\FloatTok}[1]{\textcolor[rgb]{0.77,0.35,0.07}{#1}}
\newcommand{\FunctionTok}[1]{\textcolor[rgb]{0.44,0.19,0.63}{#1}}
\newcommand{\ImportTok}[1]{\textcolor[rgb]{0.44,0.19,0.63}{#1}}
\newcommand{\InformationTok}[1]{\textcolor[rgb]{0.38,0.63,0.69}{#1}}
\newcommand{\KeywordTok}[1]{\textcolor[rgb]{0.12,0.31,0.71}{#1}}
\newcommand{\NormalTok}[1]{#1}
\newcommand{\OperatorTok}[1]{\textcolor[rgb]{0.20,0.20,0.20}{#1}}
\newcommand{\OtherTok}[1]{\textcolor[rgb]{0.00,0.44,0.13}{#1}}
\newcommand{\PreprocessorTok}[1]{\textcolor[rgb]{0.44,0.19,0.63}{#1}}
\newcommand{\RegionMarkerTok}[1]{#1}
\newcommand{\SpecialCharTok}[1]{\textcolor[rgb]{0.63,0.19,0.43}{#1}}
\newcommand{\SpecialStringTok}[1]{\textcolor[rgb]{0.63,0.19,0.43}{#1}}
\newcommand{\StringTok}[1]{\textcolor[rgb]{0.63,0.19,0.43}{#1}}
\newcommand{\VariableTok}[1]{\textcolor[rgb]{0.10,0.09,0.49}{#1}}
\newcommand{\VerbatimStringTok}[1]{\textcolor[rgb]{0.63,0.19,0.43}{#1}}
\newcommand{\WarningTok}[1]{\textcolor[rgb]{0.38,0.63,0.69}{#1}}
\makeatletter
\@ifundefined{Shaded}{}{\renewenvironment{Shaded}{\def\FrameCommand{\fboxsep=5pt\fboxrule=0.5pt\fcolorbox{wdblue!35}{codebg}}%
   \MakeFramed{\advance\hsize-\width\FrameRestore}\small}{\endMakeFramed}}
\makeatother
\fvset{fontsize=\small}

% ---------- 表格 ----------
\renewcommand{\arraystretch}{1.42}
\newcolumntype{Y}{>{\raggedright\arraybackslash}X}

% ---------- 图片 ----------
% 课件截图做成浮动体：LaTeX 会自动用正文填满页面，避免大块空白
\setcounter{topnumber}{3}\setcounter{bottomnumber}{2}\setcounter{totalnumber}{4}
\renewcommand{\topfraction}{0.85}\renewcommand{\bottomfraction}{0.6}\renewcommand{\textfraction}{0.12}\renewcommand{\floatpagefraction}{0.75}
\setlength{\textfloatsep}{10pt plus 3pt minus 2pt}\setlength{\floatsep}{8pt plus 2pt}
\newcommand{\slidefig}[3]{\begin{figure}[!ht]\centering\begin{tcolorbox}[enhanced,width=#2\linewidth,arc=2pt,boxrule=0.6pt,colframe=wdgray!70,
   colback=white,left=2pt,right=2pt,top=2pt,bottom=2pt,drop fuzzy shadow=black!25]\includegraphics[width=\linewidth]{#1}\end{tcolorbox}
   \vspace{-3pt}{\small\color{wdgray}#3}\end{figure}}
\newcommand{\plainfig}[3]{\par\noindent\begin{minipage}{\linewidth}\centering\includegraphics[width=#2\linewidth]{#1}\par\vspace{2pt}{\small\color{wdgray}#3}\end{minipage}\par\vspace{4pt}}

\begin{document}
% ---------- 封面条（仿课件封面） ----------
\thispagestyle{plain}
\begin{tcolorbox}[enhanced,enlarge left by=-4.8cm,width=\linewidth+4.8cm,arc=0pt,boxrule=0pt,
   colback=wddark,interior style={left color=wddark,right color=wdblue},left=18pt,right=18pt,top=14pt,bottom=12pt,
   overlay={\draw[white,line width=2.5pt] ([xshift=10pt,yshift=-10pt]frame.north west) -- ++(0,-0.9cm);
            \draw[white,line width=2.5pt] ([xshift=10pt,yshift=-10pt]frame.north west) -- ++(1.6cm,0);
            \draw[white,line width=2.5pt] ([xshift=-10pt,yshift=10pt]frame.south east) -- ++(0,0.9cm);
            \draw[white,line width=2.5pt] ([xshift=-10pt,yshift=10pt]frame.south east) -- ++(-1.6cm,0);}]
  {\color{white!85!wdlightblue}\small 数据结构 · 第3章
栈、队列和数组 · 康奈尔听课笔记}\par\vspace{6pt}
  {\color{white}\fontsize{22}{28}\selectfont\bfseries 3.3
课后习题：栈和队列的应用}\par\vspace{6pt}
  {\color{white!80!wdlightblue}\small 王道书 §3.3.6 试题精选 · §3.3.7
答案与解析（第 94--100 页）· 习题讲解字幕 P30--P33（选择 1--20）}
\end{tcolorbox}
\vspace{4pt}
\begin{tcolorbox}[enhanced,arc=3pt,boxrule=0.8pt,colframe=wdblue,colback=wdblue!3,left=10pt,right=10pt,top=6pt,bottom=6pt,
   colbacktitle=wdblue,coltitle=white,fonttitle=\bfseries\small,title={目录 · 点击标题跳转}]
\hypertarget{toc}{}\makeatletter\@starttoc{toc}\makeatother
\end{tcolorbox}
\hypertarget{ux8fd9ux7ec4ux9898ux5728ux8003ux4ec0ux4e48}{%
\section{这组题在考什么}\label{ux8fd9ux7ec4ux9898ux5728ux8003ux4ec0ux4e48}}

\tocback

\Needspace{4\baselineskip}

\cuenote{20 道题分成哪几类？哪几类要手动模拟，模拟出来的结果为什么算数？}20
道选择题可以分成五类：

\begin{itemize}
\tightlist
\item
  \textbf{用栈还是用队列}（01、03、09、10、12）：哪些应用用栈，哪些用队列；
\item
  \textbf{表达式}（02、19、04、13、14、18）：中缀转后缀、两个栈求中缀表达式的值、转换过程中栈里有什么、栈最深多少；
\item
  \textbf{递归}（05--08、15）：递归函数的值、调用次数、调用次序、效率、递归工作栈里的内容；
\item
  \textbf{关于栈的判断}（11、17）：两道“下列说法哪个对/错”的综合题；
\item
  \textbf{括号匹配和火车轨道}（20、16）：栈的容量够不够，队列至少要几个。
\end{itemize}

这一节的题大多要\hlred{手动模拟}。模拟为什么能当证明用，下面会反复用到两条道理：

\begin{enumerate}
\def\labelenumi{\arabic{enumi}.}
\tightlist
\item
  转换、求值这些算法每一步怎么做是规则定死的（确定性的）。对一个具体的式子从头到尾按规则走一遍，这张表就是算法的全部执行过程，\uln{不是举例，而是算法本身}，表里的结果就是答案；
\item
  判断“一定、必须、只要\ldots\ldots 就”这种全称说法时，\uln{一个反例就足以证明它错}；要证明“不一定”，找到一个成立的例子也就够了。但要证明一个全称说法对，例子再多也不够，要讲出道理（11、16、17、20
  都会用到）。
\end{enumerate}

{\rowcolors{2}{wdblue!5}{white}\arrayrulecolor{wdblue!35}
\begin{xltabular}{\linewidth}{>{\hsize=0.474\hsize\raggedright\arraybackslash}X>{\hsize=2.053\hsize\raggedright\arraybackslash}X>{\hsize=0.474\hsize\raggedright\arraybackslash}X}
\rowcolor{wdblue}\textcolor{white}{\bfseries 题号} & \textcolor{white}{\bfseries 考点} & \textcolor{white}{\bfseries 答案} \\
\endfirsthead
01 & 栈的应用不包括 & D \\ \hline
02 & \texttt{a*(b+c)-d} 的后缀表达式 & B \\ \hline
03 & 用到了队列的 & D \\ \hline
04 & 运算数栈只有两个单元，不会溢出的表达式 & B \\ \hline
05 & \texttt{f(f(1))} 的值 & B \\ \hline
06 & 计算 \texttt{F(8)} 调用递归函数的次数 & C \\ \hline
07 & \texttt{func(func(5))} 中第 4 个被执行的 & C \\ \hline
08 & 递归算法和非递归算法的效率 & B \\ \hline
09 & 局部变量采用什么结构存储 & C \\ \hline
10 & 需要队列作辅助存储空间的操作 & B \\ \hline
11 & 关于栈和队列，正确的是 & A \\ \hline
12 & 【2009 统考】打印数据缓冲区的逻辑结构 & B \\ \hline
13 & 【2012 统考】中缀转后缀时栈中操作符最多几个 & A \\ \hline
14 & 【2014 统考】扫描到 \texttt{f} 时栈中的元素 & B \\ \hline
15 & 【2015 统考】递归调用时栈底到栈顶的信息 & A \\ \hline
16 & 【2016 统考】火车轨道至少几条 & C \\ \hline
17 & 【2017 统考】关于栈，错误的是 & C \\ \hline
18 & 【2018 统考】两个栈，调用 3 次 \texttt{F()} 后的栈顶 & B \\ \hline
19 & 【2024 统考】\texttt{x+y*(z-u)/v} 的后缀表达式 & A \\ \hline
20 & 【2025 统考】容量为 3 的栈做括号匹配，处理不了的 & D \\ \hline
\end{xltabular}\par\vspace{4pt}}

\begin{fixbox}

\textbf{习题讲解视频的观看建议}：

\begin{itemize}
\tightlist
\item
  \textbf{P33 的 19
  最值得看}：一休老师讲中缀转后缀的第二种方法“搭积木”（画表达式树再后序遍历），还讲了判断优先级不用死记、用小学数学想“先算谁”。正课没有这个方法；
\item
  \textbf{P31 的 08
  值得看}：老师把斐波那契的递归和循环写法放在一起比较时间、空间复杂度，重点讲了\uln{递归的空间复杂度是 \(O(n)\) 不是 \(O(2^n)\)}，画了
  F(3) 的栈。下面 08 题我补了证明；
\item
  \textbf{P30 的 02、04}：02 的“从选项倒推”很省时间；04
  老师在板书上用两个栈完整模拟了选项 A、B；
\item
  \textbf{P33 的 20}：只看括号、用计数器算栈深的技巧，比较短；
\item
  \textbf{P32 的
  13、14}：两道统考题的完整模拟，下面都有完整的表，可以不看；
\item
  \textbf{P30 的 01、03、05，P31 的 06、07、09、10，P32 的
  11、12、15，P33 的 16、17、18}：看书上解析和下面的笔记即可。16
  题老师说“没有什么严格上的理论证明”，下面我补了证明。
\end{itemize}

下面每题的视频框都记了老师的做法。

\end{fixbox}

\hypertarget{ux4e00ux7528ux6808ux8fd8ux662fux7528ux961fux52170103091012}{%
\section{一、用栈还是用队列（01、03、09、10、12）}\label{ux4e00ux7528ux6808ux8fd8ux662fux7528ux961fux52170103091012}}

\tocback

\Needspace{4\baselineskip}

\cuenote{怎样一眼判断一个应用用栈还是用队列？}判断的依据只有一条：这个应用要求\uln{后来的先处理（后进先出）还是先来的先处理（先进先出）}。

\begin{itemize}
\tightlist
\item
  后进先出 →
  栈：括号匹配（最后出现的左括号最先被匹配）、表达式求值（暂时不能确定运算次序的运算符/操作数）、递归和函数调用（最后被调用的函数最先结束）；
\item
  先进先出 → 队列：缓冲区、FIFO
  页面替换、树的层次遍历、图的广度优先遍历、操作系统的先来先服务。
\end{itemize}

\begin{questionbox}[01]

栈的应用不包括（　）。

A. 递归　　　B. 表达式求值　　　C. 括号匹配　　　D. 缓冲区

\end{questionbox}

\hlred{\textbf{答案：D}}

\begin{videobox}

习题讲解字幕
P30：A，递归要用到函数栈，递归的过程不断把函数信息压栈，如果忘了写递归出口，栈会被压满溢出；B、C
都是正课讲过的栈的应用。D，缓冲区是操作系统里的概念，打印机等资源有限，要用的任务排队，先来的先用，所以缓冲区是队列。

\end{videobox}

\begin{jiexibox}

王道书第 96 页：缓冲区是用队列实现的，选项 A、B、C 都是栈的典型应用。

\end{jiexibox}

\begin{recallbox}

正课 3.3.1、3.3.2、3.3.3
分别是栈在括号匹配、表达式求值、递归中的应用；缓冲区在正课
3.3.4+3.3.5“三、队列在操作系统中的应用”。

\end{recallbox}

\begin{questionbox}[03]

下面（　）用到了队列。

A. 括号匹配　　　B. 表达式求值　　　C. 递归　　　D. FIFO 页面替换算法

\end{questionbox}

\hlred{\textbf{答案：D}}

\begin{videobox}

习题讲解字幕 P30：FIFO 是 First In First
Out，先进先出，从名字就能看出用的是队列。FIFO
页面替换算法在操作系统里会学；考自命题的同学一般不会考和操作系统强相关的概念。

\end{videobox}

\begin{jiexibox}

王道书第 97 页：FIFO 页面替换算法用到了队列。其余的都只用到了栈。

\end{jiexibox}

\begin{questionbox}[09]

执行函数时，其局部变量一般采用（　）进行存储。

A. 树形结构　　　B. 静态链表　　　C. 栈结构　　　D. 队列结构

\end{questionbox}

\hlred{\textbf{答案：C}}

\begin{videobox}

习题讲解字幕
P31：严格来说这是操作系统的知识。程序运行时内存里有一块区域专门为它服务，分成堆区、栈区等，局部变量一般存在栈区，所以是栈结构。

\end{videobox}

\begin{jiexibox}

王道书第 97
页：调用函数时，系统会为调用者构造一个由参数表和返回地址组成的活动记录，并将记录压入系统提供的栈中，若被调用函数有局部变量，也要压入栈中。

\end{jiexibox}

\begin{defbox}[：名词统一]

书上“活动记录（参数表和返回地址）“＋”局部变量”＝正课 3.3.3
函数调用栈保存的三样东西：①调用返回地址 ②实参
③局部变量。书上“系统提供的栈”＝正课“函数调用栈”＝老师说的“函数栈”“栈区”。

\end{defbox}

\begin{questionbox}[10]

执行（　）操作时，需要使用队列作为辅助存储空间。

A. 查找散列（哈希）表　　　　B. 广度优先搜索图\\
C. 前序（根）遍历二叉树　　　D. 深度优先搜索图

\end{questionbox}

\hlred{\textbf{答案：B}}

\begin{videobox}

习题讲解字幕
P31：这题涉及后面的知识，学完再回来看。A，散列表用到的是数组，或者邻接表那样的结构；B，图的广度优先搜索确实要用队列；C，408
只考虑递归遍历二叉树，不考虑非递归写法，递归用的是栈；D，深度优先搜索也用栈。

\end{videobox}

\begin{jiexibox}

王道书第 97 页：本题涉及第 5 章和第 6
章的内容，图的广度优先搜索类似于树的层序遍历，都要借助于队列。

\end{jiexibox}

\begin{defbox}[：名词统一]

书上“树的层序遍历”＝正课
3.3.4+3.3.5“树的层次遍历”；书上“图的广度优先搜索”＝正课“图的广度优先遍历”。

\end{defbox}

\begin{recallbox}

正课 3.3.4+3.3.5
第一、二节：层次遍历和广度优先遍历都是“处理队头、把它的孩子（或相邻且未访问的结点）加入队尾”。

\end{recallbox}

\begin{questionbox}[12]

【2009
统考真题】为解决计算机主机与打印机之间速度不匹配的问题，通常设置一个打印数据缓冲区，主机将要输出的数据依次写入该缓冲区，而打印机则依次从该缓冲区中取出数据。该缓冲区的逻辑结构应该是（　）。

A. 栈　　　B. 队列　　　C. 树　　　D. 图

\end{questionbox}

\hlred{\textbf{答案：B}}

\begin{videobox}

习题讲解字幕 P32：老师举例，主机要打印 W、D、N、B
四个字，先交给缓冲区的是 W；打印机打出来也必须先是
W，不然顺序就乱了。先进来的先出去，所以是队列。

\end{videobox}

\begin{jiexibox}

王道书第 98
页：在提取数据时必须保持原来数据的顺序，所以缓冲区的特性是先进先出。

\end{jiexibox}

\begin{recallbox}

正课 3.3.4+3.3.5 第三节的书框：王道书第 94
页考点追踪“缓冲区的逻辑结构（2009）”，就是这道题，\uln{打印数据缓冲区实际上就是一个队列}。

\end{recallbox}

\hypertarget{ux4e8cux8868ux8fbeux5f0f021904131418}{%
\section{二、表达式（02、19、04、13、14、18）}\label{ux4e8cux8868ux8fbeux5f0f021904131418}}

\tocback

\Needspace{4\baselineskip}

\cuenote{求后缀表达式有哪几种方法？每种方法凭什么是对的？}这组题用到正课
3.3.2 的全部内容。先把几个方法和它们的道理放在一起：

\begin{enumerate}
\def\labelenumi{\arabic{enumi}.}
\item
  \textbf{按运算次序组合}（正课 3.3.2\_1
  第二节，手算）：确定各运算符的运算次序，每次按“左操作数 右操作数
  运算符”组合成一个新的操作数。把这条规则写成递归形式就是：

  \begin{itemize}
  \tightlist
  \item
    单个操作数的后缀就是它自己；
  \item
    表达式 \(L \text{ op } R\)（op 是\uln{最后一个执行的运算符}）的后缀
    \(=\) \(L\) 的后缀，接 \(R\) 的后缀，再接 op。
  \end{itemize}

  最后执行的运算符怎么找：去掉所有括号里面的部分，在括号外面的运算符里找\uln{优先级最低的}；同一优先级有好几个时，因为同级从左往右算，\uln{最右边的那个最后算}。
\item
  \textbf{用栈转换}（正课 3.3.2\_2，机算）：规则见正课 3.3.2\_2 第一节。
\item
  \textbf{表达式树}（习题讲解视频 P33 的
  19，一休老师）：画出表达式树，后序遍历。其实和方法 1 是同一件事：方法
  1 每次拆出来的“最后执行的运算符”就是树根，\(L\)、\(R\)
  就是左右子树，“\(L\) 后缀、\(R\)
  后缀、op”正是后序遍历“左、右、根”的顺序。
\item
  \textbf{从选项倒推}（习题讲解视频 P30 的
  02）：把选项的后缀表达式转回中缀（正课 3.3.2\_1
  我补的“后缀转中缀”），看和题目是不是同一个式子。
\end{enumerate}

因为运算次序由优先级、括号、左结合唯一确定，方法 1
每一步的拆法都是唯一的，所以\uln{一个中缀表达式的后缀表达式是唯一的}（按左优先原则）。这一点保证了排除法是可靠的：只要选项里有一处和正确答案不符，它就错了。

\begin{questionbox}[02]

表达式 \texttt{a*(b+c)-d} 的后缀表达式是（　）。

A. \texttt{abcd*+-}　　　B. \texttt{abc+*d-}　　　C.
\texttt{abc*+d-}　　　D. \texttt{-+*abcd}

\end{questionbox}

\hlred{\textbf{答案：B}}

\begin{videobox}

习题讲解字幕
P30：两种方法。第一种，用栈模拟中缀转后缀，老师在板书上从头到尾模拟了一遍；第二种，做选择题更快：从选项倒推，把后缀转回中缀。A
里有 \texttt{cd*}，即 \texttt{c*d}，原式中没有，排除；B：\texttt{bc+} 得
\texttt{b+c}，再和 \texttt{a} 做 \texttt{*} 得 \texttt{a*(b+c)}，再和
\texttt{d} 做 \texttt{-}，正好是原式；C 里有 \texttt{bc*}，原式中没有；D
是前缀的样子，不是后缀。

\end{videobox}

\begin{jiexibox}

王道书第 97
页：后缀表达式中，每个运算符均直接位于其两个操作数的后面，按照这样的方式逐步根据计算的优先级将每个计算式进行变换，即可得到后缀表达式。

【另解】将两个直接操作数用括号括起来，再将操作符提到括号后，最后去掉括号。例如，对于
\texttt{((a*(b+c))-d)}，提出操作符并去掉括号后，可得后缀表达式为
\texttt{abc+*d-}。

学完第 5
章后，可将表达式画成二叉树的形式，再用后序遍历即可求得后缀表达式。

\end{jiexibox}

\begin{addbox}

\textbf{用上面方法 1 严格推一遍}：

\begin{enumerate}
\def\labelenumi{\arabic{enumi}.}
\tightlist
\item
  \texttt{a*(b+c)-d}：括号外的运算符是
  \texttt{*}、\texttt{-}，\texttt{-} 优先级最低，最后执行。拆成 \(L=\)
  \texttt{a*(b+c)}，\(R=\) \texttt{d}，后缀 \(=\) (\(L\) 的后缀)
  \texttt{d} \texttt{-}；
\item
  \texttt{a*(b+c)}：括号外只有 \texttt{*}。拆成 \(L=\)
  \texttt{a}，\(R=\) \texttt{b+c}，后缀 \(=\) \texttt{a} (\(R\) 的后缀)
  \texttt{*}；
\item
  \texttt{b+c} 的后缀是 \texttt{bc+}。
\end{enumerate}

逐层代回：\texttt{a} \texttt{bc+} \texttt{*} → \texttt{abc+*}；再接
\texttt{d-} → \texttt{abc+*d-}，就是 B。由上面“后缀唯一”，A、C、D
都不可能对。

老师的倒推法为什么也算证明：B 转回中缀得到的是
\texttt{(a*(b+c))-d}，和原式的运算次序完全相同，所以 B
就是原式的后缀表达式；这一步本身就证明了 B 对，不只是排除了别的选项。

\end{addbox}

\begin{recallbox}

正课 3.3.2\_1
第二节中缀转后缀（手算）的“左优先原则”，以及我补的“后缀转中缀”。

\end{recallbox}

\begin{questionbox}[19]

【2024 统考真题】与表达式 \texttt{x+y*(z-u)/v}
等价的后缀表达式是（　）。

A. \texttt{xyzu-*v/+}　　　B. \texttt{xyzu-v/*+}　　　C.
\texttt{-+x/*y-zuv}　　　D. \texttt{+x*y/-zuv}

\end{questionbox}

\hlred{\textbf{答案：A}}

\begin{videobox}

习题讲解字幕 P33（一休老师）：两种方法。

\textbf{方法一，模拟}，大家都要掌握：从左到右扫描，操作数直接加入后缀表达式；遇到运算符和栈顶比较优先级，栈顶优先级大于等于当前运算符就出栈，括号不写入后缀表达式；扫描完把栈里剩下的依次出栈。比较优先级时\uln{不用背}，用小学数学想“先算谁”：\texttt{x+y*…}
中应该先算后面的乘法，所以 \texttt{*} 入栈；括号里的要先算，所以
\texttt{(} 入栈；\texttt{-} 前面是括号，要先算括号里的
\texttt{z-u}，所以 \texttt{-} 入栈；遇到
\texttt{)}，括号里的算完，依次出栈到 \texttt{(}；遇到
\texttt{/}，\texttt{y*(z-u)} 和后面的 \texttt{/v} 应该先算左边的乘，所以
\texttt{*} 出栈、\texttt{/} 入栈；最后依次出栈
\texttt{/}、\texttt{+}。结果 \texttt{xyzu-*v/+}。

\textbf{方法二，搭积木}（求表达式树）：中缀表达式是表达式树的中序遍历，后缀表达式是后序遍历。把整个表达式
\(F\) 看成大积木，拆成左边的小积木 \(F_1\)、中间的运算符、右边的小积木
\(F_2\)，再分别递归地拆。拆的时候\hlred{必须遵守四则运算规则}，比如把
\texttt{x+y} 当成 \(F_1\) 就错了，会先算
\texttt{x+y}，“你的小学老师可能就会来找你了”。有时有多种合法拆法：\texttt{y*(z-u)/v}
可以拆成 \texttt{y} 和 \texttt{(z-u)/v}，也可以拆成 \texttt{y*(z-u)} 和
\texttt{v}。为了让表达式树唯一，规定\uln{在合法的前提下选最长的 \(F_1\)}，即
\(F_1=\) \texttt{y*(z-u)}、\(F_2=\)
\texttt{v}。拆完后序遍历（左右根），答案也是 A。

老师总结：方法一最实用，一定要掌握；方法二可以加入自己的武器库。

\end{videobox}

\begin{jiexibox}

王道书第 99
页：根据中缀表达式画出对应的二叉树如下图所示，对该二叉树进行后序遍历即可得到后缀表达式为
\texttt{xyzu-*v/+}。本题也可采用本书前面描述的手算方法。

\end{jiexibox}

\slidefig{assets/3.3ex-19-book.png}{0.22}{王道书第 99 页：19 题的表达式树}

\begin{addbox}

\textbf{用方法 1 严格推一遍，并说明“\(F_1\) 最长”的道理}：

\begin{enumerate}
\def\labelenumi{\arabic{enumi}.}
\tightlist
\item
  \texttt{x+y*(z-u)/v}：括号外的运算符是
  \texttt{+}、\texttt{*}、\texttt{/}，\texttt{+} 最低，最后执行。\(L=\)
  \texttt{x}，\(R=\) \texttt{y*(z-u)/v}；
\item
  \texttt{y*(z-u)/v}：括号外是
  \texttt{*}、\texttt{/}，同一优先级，同级从左往右算，\uln{最右边的 \texttt{/} 最后执行}。\(L=\)
  \texttt{y*(z-u)}，\(R=\) \texttt{v}；
\item
  \texttt{y*(z-u)}：\(L=\) \texttt{y}，\(R=\) \texttt{z-u}，后缀
  \texttt{yzu-*}；
\item
  代回：\texttt{yzu-*} 接 \texttt{v} 再接 \texttt{/} →
  \texttt{yzu-*v/}；前面接 \texttt{x}、最后接 \texttt{+} →
  \texttt{xyzu-*v/+}，就是 A。
\end{enumerate}

第 2 步取最右边的 \texttt{/}
作为拆分点，左边那块自然就最长。所以老师的“在合法的前提下选最长的
\(F_1\)“不是人为加的规定，它就是”同级运算从左往右算”（左结合）这条四则运算规则。

\textbf{B 为什么不算“等价”}：B 转回中缀是
\texttt{x+y*((z-u)/v)}，先除后乘，\uln{运算次序和原式不同}。在实数运算里两者数值相等，但题目说的等价指的是同一个运算过程；而且在
C 语言的整数除法下数值也可能不同，例如 \(y=2\)，\(z-u=1\)，\(v=2\)：原式
\((2\times1)/2=1\)，B 是 \(2\times(1/2)=2\times0=0\)。C、D
是前缀的样子。

\end{addbox}

\begin{recallbox}

正课 3.3.2\_1
第二节“左优先原则”：只要左边的运算符能先计算，就优先算左边的，保证手算结果和机算一致。本题第
2 步就是这条原则。

\end{recallbox}

\Needspace{8\baselineskip}\vspace{22pt}

\cueblock{\cuenote}{\relax{两个栈求中缀表达式的值时，运算数栈最深有多深？怎样不模拟也能算出来？}}

\begin{questionbox}[04]

利用栈求表达式的值时，设立运算数栈 OPEN。假设 OPEN
只有两个存储单元，则在下列表达式中，不会发生溢出的是（　）。

A. \texttt{A-B*(C-D)}　　B. \texttt{(A-B)*C-D}　　C.
\texttt{(A-B*C)-D}　　D. \texttt{(A-B)*(C-D)}

\end{questionbox}

\hlred{\textbf{答案：B}}

\begin{videobox}

习题讲解字幕
P30：这是中缀表达式求值。理论上可以分两步：先中缀转后缀，再后缀求值；咸鱼老师的基础课讲过这两步可以合二为一，用一个运算符栈、一个运算数栈。

\begin{itemize}
\tightlist
\item
  \textbf{A}：\texttt{A} 入运算数栈；\texttt{-} 入运算符栈；\texttt{B}
  入栈；\texttt{*} 比 \texttt{-} 高，入栈；\texttt{(} 入栈；扫描到
  \texttt{C} 时运算数栈已经有 \texttt{A}、\texttt{B}
  两个，满了，\texttt{C} 无法入栈，溢出；
\item
  \textbf{B}：\texttt{(} 入栈；\texttt{A} 入栈；\texttt{-}
  入栈；\texttt{B} 入栈；遇到 \texttt{)}，\texttt{-}
  出栈，同时运算数栈弹出两个算 \texttt{A-B}，结果压回；\texttt{(} 和
  \texttt{)} 抵消；\texttt{*} 入栈；\texttt{C} 入栈；遇到 \texttt{-}，比
  \texttt{*} 低，\texttt{*} 出栈，算 \texttt{(A-B)*C} 压回；\texttt{-}
  入栈；\texttt{D} 入栈；最后 \texttt{-} 出栈，算完。整个过程没有溢出；
\item
  \textbf{C、D}：老师让同学按 A、B
  的过程自己模拟，都会溢出。他又用另一种方法做了 C：先转后缀
  \texttt{ABC*-D-}，后缀求值时 \texttt{A}、\texttt{B}、\texttt{C}
  都要先入栈，两个单元不够，也会溢出。老师强调：中缀求值只是把两个过程合二为一，\uln{没有本质上的改变}，所以后缀求值时运算数栈不够，用哪种方法算都会溢出。D
  让同学自己验证。
\end{itemize}

\end{videobox}

\begin{jiexibox}

王道书第 97
页：利用栈求表达式的值时，可以分别设立运算符栈和运算数栈，其原理不变。选项
B 中 A 入栈，B 入栈，计算得 R1，C 入栈，计算得 R2，D 入栈，计算得
R3，由此得栈深为 2。选项 A、C、D 依次计算得栈深为 4、3、3。因此选择选项
B。

\textbf{技巧}：根据运算符优先级，统计已依次入栈但还未参与运算的运算符数。以选项
C 为例，当“(“、”A”、“-”入栈时，“(”和“-“还未参与运算，此时运算符栈大小为
2，”B”和“*“入栈时运算符栈大小为 3，”C”入栈时“B*C”运算，运算符栈大小为
2，以此类推。

\end{jiexibox}

\begin{addbox}

\textbf{补全老师留给同学的 C、D 模拟}（双栈，左边是栈底）：

{\rowcolors{2}{wdblue!5}{white}\arrayrulecolor{wdblue!35}
\begin{xltabular}{\linewidth}{>{\hsize=0.400\hsize\raggedright\arraybackslash}X>{\hsize=0.880\hsize\raggedright\arraybackslash}X>{\hsize=1.200\hsize\raggedright\arraybackslash}X>{\hsize=1.520\hsize\raggedright\arraybackslash}X}
\rowcolor{wdblue}\textcolor{white}{\bfseries 扫描} & \textcolor{white}{\bfseries 运算符栈} & \textcolor{white}{\bfseries 运算数栈} & \textcolor{white}{\bfseries 动作} \\
\endfirsthead
\textbf{C} &  &  & \texttt{(A-B*C)-D} \\ \hline
( A & ( & A &  \\ \hline
- B & (- & A, B & 栈顶是 \texttt{(}，\texttt{-} 入栈 \\ \hline
* & (-* & A, B & \texttt{*} 比 \texttt{-} 高，入栈 \\ \hline
C & (-* & A, B, \hlred{C} & \hlred{第 3 个，溢出} \\ \hline
\textbf{D} &  &  & \texttt{(A-B)*(C-D)} \\ \hline
( A - B & (- & A, B &  \\ \hline
) & 空 & R1 & \texttt{-} 出栈，算 R1＝A−B；丢弃 \texttt{(} \\ \hline
* ( & *( & R1 &  \\ \hline
C & *( & R1, C &  \\ \hline
- & *(- & R1, C & 栈顶是 \texttt{(}，入栈 \\ \hline
D & *(- & R1, C, \hlred{D} & \hlred{第 3 个，溢出} \\ \hline
\end{xltabular}\par\vspace{4pt}}

\textbf{不模拟，直接算出栈深（证明）}：

\begin{enumerate}
\def\labelenumi{\arabic{enumi}.}
\item
  \emph{双栈算法的运算数栈和后缀求值的栈是同一个栈。} 正课 3.3.2\_2
  第五节说双栈算法是“一边生成后缀表达式，一边计算”：转换算法每输出一个操作数，双栈算法就把它压入运算数栈；转换算法每弹出（输出）一个运算符，双栈算法就弹出两个操作数、算完压回。这正好是后缀求值对同一个记号做的动作。所以每输出一个记号后，两边的栈内容完全相同，\uln{最大深度也相同}。这就是书上“其原理不变”、老师“没有本质上的改变”的意思。
\item
  \emph{后缀求值的栈深公式。} 后缀求值时，读到操作数压入 1 个，栈深
  \(+1\)；读到（双目）运算符弹出 2 个、压回 1 个，栈深
  \(-1\)。所以读完前 \(k\) 个记号后，

  \[\text{栈深} = (\text{前 } k \text{ 个记号中的操作数个数}) - (\text{其中的运算符个数}).\]

  对所有 \(k\) 取最大值，就是运算数栈的最大深度。
\item
  \emph{代入四个选项}（先按方法 1 写出后缀，再逐个记号写出栈深）：
\end{enumerate}

{\rowcolors{2}{wdblue!5}{white}\arrayrulecolor{wdblue!35}
\begin{xltabular}{\linewidth}{>{\hsize=0.880\hsize\raggedright\arraybackslash}X>{\hsize=0.880\hsize\raggedright\arraybackslash}X>{\hsize=1.440\hsize\raggedright\arraybackslash}X>{\hsize=0.800\hsize\raggedright\arraybackslash}X}
\rowcolor{wdblue}\textcolor{white}{\bfseries 选项} & \textcolor{white}{\bfseries 后缀表达式} & \textcolor{white}{\bfseries 每读一个记号后的栈深} & \textcolor{white}{\bfseries 最大} \\
\endfirsthead
A \texttt{A-B*(C-D)} & \texttt{ABCD-*-} & 1, 2, 3, 4, 3, 2, 1 & 4 \\ \hline
B \texttt{(A-B)*C-D} & \texttt{AB-C*D-} & 1, 2, 1, 2, 1, 2, 1 & \hlred{2} \\ \hline
C \texttt{(A-B*C)-D} & \texttt{ABC*-D-} & 1, 2, 3, 2, 1, 2, 1 & 3 \\ \hline
D \texttt{(A-B)*(C-D)} & \texttt{AB-CD-*} & 1, 2, 1, 2, 3, 2, 1 & 3 \\ \hline
\end{xltabular}\par\vspace{4pt}}

只有 B 的最大深度不超过 2，和书上“栈深 4、3、3”一致。

\end{addbox}

\begin{fixbox}

\textbf{书上的“技巧”要小心括号}。书上说统计“已入栈但还未参与运算的运算符数”，可它的例子把
\texttt{(}
也数进了“运算符栈大小”。准确的说法是：\uln{每读入一个操作数之后，运算数栈中的个数＝运算符栈中（不算括号的）运算符个数＋1}。

\emph{证明}：记 \(D=\) 运算数个数 \(-\)
运算符栈中非括号运算符个数。读入操作数，\(D+1\)；运算符入栈，\(D-1\)；弹出一个运算符并计算，运算数
\(-1\)（弹 2 压 1）、运算符 \(-1\)，\(D\) 不变；括号入栈、出栈都不影响
\(D\)。合法的中缀表达式去掉括号后是“操作数、运算符、操作数、\ldots\ldots、操作数”交替出现，开始时
\(D=0\)，所以每读完一个操作数 \(D=1\)，每压入一个运算符后
\(D=0\)。又因为运算数栈只在读入操作数时变多，所以最大深度一定出现在刚读完某个操作数时，等于那时\uln{未参与运算的运算符数＋1}。

\emph{例子}：选项 C 读入 \texttt{C} 时，运算符栈是
\texttt{(\ -\ *}，不算括号有 2 个运算符，运算数栈有 \(2+1=3\)
个。书上数成“运算符栈大小为 3”是因为恰好只有一个
\texttt{(}，碰巧相等。换成 \texttt{((A-B))}：读入 \texttt{B}
时运算符栈是 \texttt{(\ (\ -}，大小为 3，可运算数栈只有 \(1+1=2\) 个。

\end{fixbox}

\begin{recallbox}

正课 3.3.2\_2
第五节“中缀表达式的计算（用栈实现）“：每弹出一个运算符，就弹出两个操作数运算，\uln{先弹出的是右操作数}。正课
3.3.2\_1 第四节”后缀表达式的计算（机算）”。

\end{recallbox}

\Needspace{8\baselineskip}\vspace{22pt}

\cueblock{\cuenote}{\relax{中缀转后缀时栈里存什么？问“栈中最多几个”“扫描到某处时栈里有什么”怎么做？}}

\begin{questionbox}[13]

【2012 统考真题】已知操作符包括“+”“-”“*”“/”“(”和“)”。将中缀表达式
\texttt{a+b-a*((c+d)/e-f)+g} 转换为等价的后缀表达式
\texttt{ab+acd+e/f-*-g+}
时，用栈来存放暂时还不能确定运算次序的操作符。栈初始时为空，转换过程中同时保存在栈中的操作符的最大个数是（　）。

A. 5　　　B. 7　　　C. 8　　　D. 11

\end{questionbox}

\hlred{\textbf{答案：A}}

\begin{videobox}

习题讲解字幕
P32：按中缀转后缀的规则完整模拟一遍，边模拟边记栈里的个数。扫描到
\texttt{d} 前后，栈里是 \texttt{-\ *\ (\ (\ +}，共 5
个，这是最多的时候；遇到后面的 \texttt{)} 才开始弹出。答案 A。

\end{videobox}

\begin{jiexibox}

王道书第 98
页：在中缀表达式转后缀表达式的过程中，扫描到操作数时直接输出，扫描到操作符时根据其优先级进行相应的出入栈操作。有几点需要注意：①
若遇到界限符“(“，则直接入栈；②
若遇到界限符”)“，则不入栈，且依次弹出栈顶运算符，直到遇到”(“为止，并删除”(“；③
若当前运算符的优先级高于栈顶运算符或遇到栈顶为”(“，则直接入栈；④
若当前运算符的优先级低于或等于栈顶运算符，则依次弹出栈中的运算符并加入后缀表达式，直到遇到一个优先级低于它的运算符或遇到”(”或栈空为止，之后将当前运算符入栈。

求栈中操作符的最大个数时，为简单起见，可省略对操作数的处理。（书上列了
12 步的表）由上述过程可知，栈中操作符的最大个数为 5。

\end{jiexibox}

\begin{defbox}[：名词统一]

书上这四条规则＝正课 3.3.2\_2 第一节的规则：书 ③④
合起来就是正课“依次弹出栈中优先级高于或等于当前运算符的所有运算符，碰到
\texttt{(}
或栈空则停止，再把当前运算符入栈”。本题的“操作符”包括括号，所以
\texttt{(} 也算个数，和正课“\texttt{(} 也占位置”一致。

\end{defbox}

\begin{addbox}

\textbf{完整的表}（书上省略了操作数，这里把后缀表达式也写出来，最后一行可以和题目给的后缀对照）：

{\rowcolors{2}{wdblue!5}{white}\arrayrulecolor{wdblue!35}
\begin{xltabular}{\linewidth}{>{\hsize=0.350\hsize\raggedright\arraybackslash}X>{\hsize=0.800\hsize\raggedright\arraybackslash}X>{\hsize=0.350\hsize\raggedright\arraybackslash}X>{\hsize=1.500\hsize\raggedright\arraybackslash}X>{\hsize=2.000\hsize\raggedright\arraybackslash}X}
\rowcolor{wdblue}\textcolor{white}{\bfseries 扫描} & \textcolor{white}{\bfseries 栈（左边是栈底）} & \textcolor{white}{\bfseries 个数} & \textcolor{white}{\bfseries 后缀表达式} & \textcolor{white}{\bfseries 动作} \\
\endfirsthead
a &  & 0 & a &  \\ \hline
+ & + & 1 & a & 栈空，入栈 \\ \hline
b & + & 1 & ab &  \\ \hline
- & - & 1 & ab+ & \texttt{+} 和 \texttt{-} 同级，弹出 \texttt{+}；\texttt{-} 入栈 \\ \hline
a & - & 1 & ab+a &  \\ \hline
* & -* & 2 & ab+a & \texttt{*} 比 \texttt{-} 高，入栈 \\ \hline
( & -*( & 3 & ab+a &  \\ \hline
( & -*(( & 4 & ab+a &  \\ \hline
c & -*(( & 4 & ab+ac &  \\ \hline
+ & \textbf{-*((+} & \hlred{5} & ab+ac & 栈顶是 \texttt{(}，入栈 \\ \hline
d & -*((+ & 5 & ab+acd &  \\ \hline
) & -*( & 3 & ab+acd+ & 弹出 \texttt{+}，删除 \texttt{(} \\ \hline
/ & -*(/ & 4 & ab+acd+ & 栈顶是 \texttt{(}，入栈 \\ \hline
e & -*(/ & 4 & ab+acd+e &  \\ \hline
- & -*(- & 4 & ab+acd+e/ & 弹出 \texttt{/}；遇 \texttt{(} 停；\texttt{-} 入栈 \\ \hline
f & -*(- & 4 & ab+acd+e/f &  \\ \hline
) & -* & 2 & ab+acd+e/f- & 弹出 \texttt{-}，删除 \texttt{(} \\ \hline
+ & + & 1 & ab+acd+e/f-*- & 依次弹出 \texttt{*}、\texttt{-}（都不低于 \texttt{+}）；\texttt{+} 入栈 \\ \hline
g & + & 1 & ab+acd+e/f-*-g &  \\ \hline
结束 &  & 0 & ab+acd+e/f-*-g+ & 弹出 \texttt{+} \\ \hline
\end{xltabular}\par\vspace{4pt}}

这张表是转换算法对这个式子的全部执行过程（每一步怎么做由规则唯一决定），所以“最多
5
个”是精确答案。最后一行和题目给的后缀表达式一致，说明每一步都没有用错规则。

\end{addbox}

\begin{recallbox}

正课 3.3.2\_2
第四节“练习：栈是否会溢出？“：老师说考试可能考\uln{转换到某一步时栈里的情况}，频率还不低；那道练习栈中最多
7 个。书上考点追踪”栈的深度分析”正好列了 2012（本题）。

\end{recallbox}

\begin{questionbox}[14]

【2014 统考真题】假设栈初始为空，将中缀表达式 \texttt{a/b+(c*d-e*f)/g}
转换为等价的后缀表达式的过程中，当扫描到 \texttt{f}
时，栈中的元素依次是（　）。

A. \texttt{+(*-}　　　B. \texttt{+(-*}　　　C. \texttt{/+(*-*}　　　D.
\texttt{/+-*}

\end{questionbox}

\hlred{\textbf{答案：B}}

\begin{videobox}

习题讲解字幕 P32：按规则模拟，扫描到 \texttt{f} 时栈里从栈底到栈顶是
\texttt{+\ (\ -\ *}，选 B；老师接着把转换做完，得到
\texttt{ab/cd*ef*-g/+}。

\end{videobox}

\begin{jiexibox}

王道书第 98--99 页：中缀表达式 \texttt{a/b+(c*d-e*f)/g}
转换为后缀表达式的过程如下（书上列了 17 步的表）。由此可知，当扫描到
\texttt{f} 时，栈中的元素依次是 \texttt{+(-*}。

【另解】采用手算方法，得出中缀式 \texttt{a/b+(c*d-e*f)/g} 对应的后缀式为
\texttt{ab/cd*ef*-g/+}。中缀表达式转后缀表达式时，操作数都直接输出，因此操作数的顺序是固定的。扫描到
\texttt{f} 时，在后缀表达式 \texttt{f}
后面的运算符要么还未入栈，要么还在栈中，需要结合中缀式来判断，\texttt{f}
后面依次出栈的运算符为 \texttt{*-/+}，\texttt{/} 在中缀表达式 \texttt{f}
之后，此时还未入栈，因此栈中的运算符（从栈底到栈顶）为
\texttt{+-*}；此外，已入栈的界限符 \texttt{(}（此时还未消解，因此
\texttt{(} 也还在栈中），栈中的元素依次是 \texttt{+(-*}。

\end{jiexibox}

\begin{addbox}

\textbf{完整的表}：

{\rowcolors{2}{wdblue!5}{white}\arrayrulecolor{wdblue!35}
\begin{xltabular}{\linewidth}{>{\hsize=0.320\hsize\raggedright\arraybackslash}X>{\hsize=0.720\hsize\raggedright\arraybackslash}X>{\hsize=1.200\hsize\raggedright\arraybackslash}X>{\hsize=1.760\hsize\raggedright\arraybackslash}X}
\rowcolor{wdblue}\textcolor{white}{\bfseries 扫描} & \textcolor{white}{\bfseries 栈（左边是栈底）} & \textcolor{white}{\bfseries 后缀表达式} & \textcolor{white}{\bfseries 动作} \\
\endfirsthead
a &  & a &  \\ \hline
/ & / & a & 栈空，入栈 \\ \hline
b & / & ab &  \\ \hline
+ & + & ab/ & \texttt{/} 比 \texttt{+} 高，弹出；\texttt{+} 入栈 \\ \hline
( & +( & ab/ &  \\ \hline
c & +( & ab/c &  \\ \hline
* & +(* & ab/c & 栈顶是 \texttt{(}，入栈 \\ \hline
d & +(* & ab/cd &  \\ \hline
- & +(- & ab/cd* & 弹出 \texttt{*}；遇 \texttt{(} 停；\texttt{-} 入栈 \\ \hline
e & +(- & ab/cd*e &  \\ \hline
* & +(-* & ab/cd*e & \texttt{*} 比 \texttt{-} 高，入栈 \\ \hline
f & \textbf{\hlred{+(-*}} & ab/cd*ef & ← 题目问的时刻 \\ \hline
) & + & ab/cd*ef*- & 依次弹出 \texttt{*}、\texttt{-}，删除 \texttt{(} \\ \hline
/ & +/ & ab/cd*ef*- & \texttt{/} 比 \texttt{+} 高，入栈 \\ \hline
g & +/ & ab/cd*ef*-g &  \\ \hline
结束 &  & ab/cd*ef*-g/+ & 依次弹出 \texttt{/}、\texttt{+} \\ \hline
\end{xltabular}\par\vspace{4pt}}

\textbf{书上另解的道理}（为什么不模拟也能定出栈里的运算符）：设扫描到操作数
\texttt{f}，它立即被输出。看中缀式里 \texttt{f} 之前的任意一个运算符
\(x\)：

\begin{itemize}
\tightlist
\item
  如果 \(x\) 在扫描到 \texttt{f} 之前已经出栈，它就在 \texttt{f}
  之前被输出，在后缀式里排在 \texttt{f} 前面；
\item
  如果 \(x\) 此时还在栈里，它以后才输出，在后缀式里排在 \texttt{f}
  后面。
\end{itemize}

中缀式里 \texttt{f}
之后的运算符此时还没扫描到，不在栈里。所以：\uln{此时栈里的运算符＝中缀式中在 \texttt{f} 之前、后缀式中在 \texttt{f}
之后的运算符}；而栈只在栈顶进出，留在栈里的元素保持入栈的先后次序，所以从栈底到栈顶就是它们在中缀式里的次序。

代入本题：中缀式中 \texttt{f} 之前的运算符依次是
\texttt{/}、\texttt{+}、\texttt{*}（第一个）、\texttt{-}、\texttt{*}（第二个）；后缀式
\texttt{ab/cd*ef*-g/+} 中 \texttt{f} 之后是
\texttt{*}（第二个）、\texttt{-}、\texttt{/}（第二个，中缀中在
\texttt{f} 之后）、\texttt{+}。两边都有的是
\texttt{+}、\texttt{-}、\texttt{*}，按中缀次序是
\texttt{+\ -\ *}。再加上还没遇到 \texttt{)} 的 \texttt{(}（它在
\texttt{+} 之后、\texttt{-} 之前入栈），得到 \texttt{+(-*}。

\end{addbox}

\begin{questionbox}[18]

【2018 统考真题】若栈 S1 中保存整数，栈 S2 中保存运算符，函数
\texttt{F()} 依次执行下述各步操作：

1）从 S1 中依次弹出两个操作数 a 和 b。 2）从 S2 中弹出一个运算符 op。
3）执行相应的运算 b op a。 4）将运算结果压入 S1 中。

假定 S1 中的操作数依次是 5, 8, 3, 2（2 在栈顶），S2 中的运算符依次是
*、-、+（+ 在栈顶）。调用 3 次 \texttt{F()} 后，S1
栈顶保存的值是（　）。

A. −15　　　B. 15　　　C. −20　　　D. 20

\end{questionbox}

\hlred{\textbf{答案：B}}

\begin{videobox}

习题讲解字幕
P33：老师说这种题是他最喜欢的，“不用动脑子，只要把自己当做一个做题的机器”，题目说什么就照着做，调用三次就做三次。结果
15，而且这时 S1 里只剩 15 一个值。

\end{videobox}

\begin{jiexibox}

王道书第 99 页：第一次调用：① 从 S1 中弹出 2 和 3；② 从 S2 中弹出 +；③
执行 3+2=5；④ 将 5 压入 S1 中，第一次调用结束后 S1 中剩余 5、8、5（5
在栈顶），S2 中剩余 *、-（- 在栈顶）。第二次调用：① 从 S1 中弹出 5 和
8；② 从 S2 中弹出 -；③ 执行 8−5=3；④ 将 3 压入 S1 中，第二次调用结束后
S1 中剩余 5、3（3 在栈顶），S2 中剩余 *。第三次调用：① 从 S1 中弹出 3 和
5；② 从 S2 中弹出 *；③ 执行 5*3=15；④ 将 15 压入 S1 中，第三次调用结束后
S1 中仅剩余 15，S2 为空。

\end{jiexibox}

\begin{addbox}

\textbf{逐次的表}（栈都是左边为栈底）：

{\rowcolors{2}{wdblue!5}{white}\arrayrulecolor{wdblue!35}
\begin{xltabular}{\linewidth}{>{\hsize=0.600\hsize\raggedright\arraybackslash}X>{\hsize=0.960\hsize\raggedright\arraybackslash}X>{\hsize=0.600\hsize\raggedright\arraybackslash}X>{\hsize=0.960\hsize\raggedright\arraybackslash}X>{\hsize=1.560\hsize\raggedright\arraybackslash}X>{\hsize=1.320\hsize\raggedright\arraybackslash}X}
\rowcolor{wdblue}\textcolor{white}{\bfseries 调用} & \textcolor{white}{\bfseries 弹出 a, b} & \textcolor{white}{\bfseries op} & \textcolor{white}{\bfseries b op a} & \textcolor{white}{\bfseries 之后的 S1} & \textcolor{white}{\bfseries 之后的 S2} \\
\endfirsthead
开始 &  &  &  & 5, 8, 3, 2 & *, -, + \\ \hline
第 1 次 & a＝2, b＝3 & + & 3+2＝5 & 5, 8, 5 & *, - \\ \hline
第 2 次 & a＝5, b＝8 & - & 8−5＝3 & 5, 3 & * \\ \hline
第 3 次 & a＝3, b＝5 & * & 5×3＝\hlred{15} & 15 & 空 \\ \hline
\end{xltabular}\par\vspace{4pt}}

最容易错的是第 2 次：先弹出的 a 是右操作数，算的是 \(8-5\) 不是
\(5-8\)，否则会得到 \(-3\)，最后得 \(-15\)（选项 A 就是这个陷阱）。

\end{addbox}

\begin{recallbox}

\texttt{F()} 就是正课 3.3.2\_2
第五节双栈求值里“弹出一个运算符，再弹出两个操作数运算，结果压回”的一步；“b
op a”就是正课 3.3.2\_1 第四节说的\uln{先出栈的是右操作数}。

\end{recallbox}

\hypertarget{ux4e09ux9012ux5f52050815}{%
\section{三、递归（05--08、15）}\label{ux4e09ux9012ux5f52050815}}

\tocback

\Needspace{4\baselineskip}

\cuenote{递归函数的调用次数、调用次序、空间复杂度，各用什么工具求？}这组题的共同工具是\textbf{递归调用树}：每一次函数调用画一个结点，它直接调用的函数画成它的孩子。

\begin{itemize}
\tightlist
\item
  调用次数＝树的结点数；
\item
  调用（开始执行）的次序＝树的先序遍历次序（07 证明）；
\item
  任一时刻递归工作栈里的内容＝从树根到当前结点的路径，所以栈的最大深度＝树的高度（08、15）。
\end{itemize}

\begin{questionbox}[05]

执行下列语句段后，i 的值为（　）。

\begin{Shaded}
\begin{Highlighting}[]
\DataTypeTok{int}\NormalTok{ f}\OperatorTok{(}\DataTypeTok{int}\NormalTok{ x}\OperatorTok{)\{}
    \ControlFlowTok{return} \OperatorTok{((}\NormalTok{x}\OperatorTok{\textgreater{}}\DecValTok{0}\OperatorTok{)?}\NormalTok{ x}\OperatorTok{*}\NormalTok{f}\OperatorTok{(}\NormalTok{x}\OperatorTok{{-}}\DecValTok{1}\OperatorTok{):}\DecValTok{2}\OperatorTok{);}
\OperatorTok{\}}
\DataTypeTok{int}\NormalTok{ i}\OperatorTok{;}
\NormalTok{i}\OperatorTok{=}\NormalTok{f}\OperatorTok{(}\NormalTok{f}\OperatorTok{(}\DecValTok{1}\OperatorTok{));}
\end{Highlighting}
\end{Shaded}

A. 2　　　B. 4　　　C. 8　　　D. 无限递归

\end{questionbox}

\hlred{\textbf{答案：B}}

\begin{videobox}

习题讲解字幕 P30：先把三元表达式翻译成 if-else：如果
\texttt{x\textgreater{}0} 成立，返回冒号前面的
\texttt{x*f(x-1)}；否则返回冒号后面的 2。先算里层
\texttt{f(1)}：\(1>0\)，返回 \(1\times f(0)\)；\(f(0)\) 中 \(0>0\)
不成立，返回 2，所以 \(f(1)=2\)。再算 \(f(2)=2\times f(1)=2\times2=4\)。

\end{videobox}

\begin{jiexibox}

王道书第 97
页：栈与递归有着紧密的联系。递归模型包括递归出口和递归体两个方面。递归出口是递归算法的出口，即终止递归的条件。递归体是一个递推的关系式。根据题意有

\texttt{f(0)=2;} \texttt{f(1)=1*f(0)=2;}
\texttt{f(f(1))=f(2)=2*f(1)=4}。

\end{jiexibox}

\begin{addbox}

\textbf{为什么不是“无限递归”}：递归出口是 \texttt{x\textless{}=0}（返回
2）。每次递归调用的参数都比当前小
1，参数是整数，从非负数出发一定会在有限步后到达 0，所以对任意 \(x\ge0\)
都会终止。本题 \texttt{f(f(1))} 的执行顺序是先算出实参
\texttt{f(1)}，再调用外层：

\begin{enumerate}
\def\labelenumi{\arabic{enumi}.}
\tightlist
\item
  调用 \texttt{f(1)} → 调用 \texttt{f(0)} → 返回 2 → \texttt{f(1)} 返回
  \(1\times2=2\)；
\item
  调用 \texttt{f(2)} → 调用 \texttt{f(1)} → 调用 \texttt{f(0)} → 返回 2
  → \texttt{f(1)} 返回 2 → \texttt{f(2)} 返回 \(2\times2=4\)。
\end{enumerate}

\end{addbox}

\begin{recallbox}

正课 3.3.3
第二节：实现递归需要\uln{递归表达式（递归体）和边界条件（递归出口）}。书上“递归模型包括递归出口和递归体”是同一件事。

\end{recallbox}

\begin{questionbox}[06]

设有如下递归函数，则计算 \texttt{F(8)}
需要调用该递归函数的次数为（　）。

\begin{Shaded}
\begin{Highlighting}[]
\DataTypeTok{int}\NormalTok{ F}\OperatorTok{(}\DataTypeTok{int}\NormalTok{ n}\OperatorTok{)\{}
    \ControlFlowTok{if}\OperatorTok{(}\NormalTok{n}\OperatorTok{\textless{}=}\DecValTok{3}\OperatorTok{)} \ControlFlowTok{return} \DecValTok{1}\OperatorTok{;}
    \ControlFlowTok{else} \ControlFlowTok{return}\NormalTok{ F}\OperatorTok{(}\NormalTok{n}\OperatorTok{{-}}\DecValTok{2}\OperatorTok{)+}\NormalTok{F}\OperatorTok{(}\NormalTok{n}\OperatorTok{{-}}\DecValTok{4}\OperatorTok{)+}\DecValTok{1}\OperatorTok{;}
\OperatorTok{\}}
\end{Highlighting}
\end{Shaded}

A. 7　　　B. 8　　　C. 9　　　D. 10

\end{questionbox}

\hlred{\textbf{答案：C}}

\begin{videobox}

习题讲解字幕 P31：把递归一层层画出来：计算 \texttt{F(8)} 要算
\texttt{F(6)}、\texttt{F(4)}；\texttt{F(6)} 要算
\texttt{F(4)}、\texttt{F(2)}；每个 \texttt{F(4)} 要算
\texttt{F(2)}、\texttt{F(0)}；\texttt{F(2)}、\texttt{F(0)}
直接返回。老师提醒：\hlred{第一次传进来的 \texttt{F(8)} 也算一次调用}，题目问的是调用该递归函数
F 的次数。一共数出 9 次。

\end{videobox}

\begin{jiexibox}

王道书第 97 页：计算 \texttt{F(8)}
的递归调用树如下图所示。由图可知，递归函数 \texttt{F()} 调用的次数为 9。

\end{jiexibox}

\slidefig{assets/3.3ex-06-book.png}{0.45}{王道书第 97 页：F(8) 的递归调用树，9 个结点}

\begin{addbox}

\textbf{用递推式严格算出次数}：记 \(C(n)\) 为计算 \texttt{F(n)} 时调用 F
的总次数（包括 \texttt{F(n)} 自己）。由函数的写法：

\begin{itemize}
\tightlist
\item
  \(n\le3\) 时，直接返回，\(C(n)=1\)；
\item
  \(n>3\)
  时，自己一次，加上两个子调用各自的全部次数，\(C(n)=1+C(n-2)+C(n-4)\)。
\end{itemize}

于是
\(C(0)=C(2)=1\)，\(C(4)=1+C(2)+C(0)=3\)，\(C(6)=1+C(4)+C(2)=5\)，\(C(8)=1+C(6)+C(4)=1+5+3=9\)。

这和数调用树的结点数是一回事：\(C(n)\) 就是以 \texttt{F(n)}
为根的子树的结点数。用递推式的好处是重复出现的子树（这里 \texttt{F(4)}
出现两次）只算一次。

\end{addbox}

\begin{questionbox}[07]

设有如下递归函数，在 \texttt{func(func(5))} 的执行过程中，第 4
个被执行的 func 函数是（　）。

\begin{Shaded}
\begin{Highlighting}[]
\DataTypeTok{int}\NormalTok{ func}\OperatorTok{(}\DataTypeTok{int}\NormalTok{ x}\OperatorTok{)\{}
    \ControlFlowTok{if}\OperatorTok{(}\NormalTok{x}\OperatorTok{\textless{}=}\DecValTok{3}\OperatorTok{)} \ControlFlowTok{return} \DecValTok{2}\OperatorTok{;}
    \ControlFlowTok{else} \ControlFlowTok{return}\NormalTok{ func}\OperatorTok{(}\NormalTok{x}\OperatorTok{{-}}\DecValTok{2}\OperatorTok{)+}\NormalTok{func}\OperatorTok{(}\NormalTok{x}\OperatorTok{{-}}\DecValTok{4}\OperatorTok{);}
\OperatorTok{\}}
\end{Highlighting}
\end{Shaded}

A. \texttt{func(2)}　　B. \texttt{func(3)}　　C. \texttt{func(4)}　　D.
\texttt{func(5)}

\end{questionbox}

\hlred{\textbf{答案：C}}

\begin{videobox}

习题讲解字幕 P31：先算里层 \texttt{func(5)}：5 大于 3，要先算
\texttt{func(3)} 再算 \texttt{func(1)}，两个都小于等于 3，各返回 2，所以
\texttt{func(5)=4}。再把 4 作为参数传进去，计算
\texttt{func(4)}。\texttt{func(5)} 是第 1 个，\texttt{func(3)} 第 2
个，\texttt{func(1)} 第 3 个，第 4 个就是 \texttt{func(4)}。

\end{videobox}

\begin{jiexibox}

王道书第 97 页：首先，执行内层参数
\texttt{func(5)=func(3)+func(1)=4}，共执行 3 次 func 函数。然后，执行
\texttt{func(func(5))=func(4)=func(2)+func(0)=4}，因此第 4 个被执行的
func 函数是
\texttt{func(4)}。也可以采用画出递归调用树的方式，即某函数的执行次序等于其在递归调用树的先序遍历中的次序。

\end{jiexibox}

\Needspace{14\baselineskip}

07 题的两棵递归调用树（我画的，圈号是开始执行的次序）：

\begin{figure}[!ht]\centering\adjustimage{scale=0.85,max width=\linewidth,max height=0.6\textheight}{/mnt/project-files/notes/1数据结构/笔记/_build/cache/cdb2582a6af8349bf4d274cc404d019512238506.pdf}\end{figure}

\begin{addbox}

\textbf{为什么执行次序就是先序遍历}（对树的高度归纳）：调用一个结点
\(v\) 时，\(v\)
先开始执行；接着它调用第一个孩子，这个孩子以及它引起的所有调用（即第一棵子树）全部执行完、返回以后，\(v\)
才调用第二个孩子。由归纳假设，每棵子树内部的次序是该子树的先序；于是整体次序是“根，第一棵子树的先序，第二棵子树的先序”，这正是先序遍历的定义。

\textbf{一个前提}：上面默认 \texttt{func(x-2)+func(x-4)} 先算左边。C
语言标准其实没有规定 \texttt{+}
两边谁先求值，考试题都按从左到右理解，书和老师也是这样做的。

\end{addbox}

\begin{questionbox}[08]

对于一个问题的递归算法求解和其对应的非递归算法求解，（　）。

A. 递归算法通常效率高一些　　　B. 非递归算法通常效率高一些\\
C. 两者相同　　　　　　　　　　D. 无法比较

\end{questionbox}

\hlred{\textbf{答案：B}}

\begin{videobox}

习题讲解字幕
P31：以斐波那契数列为例，\(F_0=0\)，\(F_1=1\)，\(F_n=F_{n-1}+F_{n-2}\)。递归写法三行：两个递归出口，再返回
\texttt{F(n-1)+F(n-2)}。非递归写法：\(F_2\) 要一轮加法，\(F_3\)
两轮，\(F_n\) 就是 \(n-1\) 轮，用
\texttt{pre1}、\texttt{pre2}、\texttt{current} 三个临时变量循环。

\begin{itemize}
\tightlist
\item
  时间：非递归 \(O(n)\)；递归像二叉树一样展开，第一层 1 次，第二层
  \(2^1\)，第三层 \(2^2\)\ldots\ldots 是指数级的
  \(O(2^n)\)，老师说可以当结论记；
\item
  空间：非递归只用几个临时变量，\(O(1)\)；递归要用函数栈，空间取决于\hlred{函数栈最大会使用多少空间，而不是一共使用了多少空间}。很多同学误以为也是
  \(O(2^n)\)，这是错的：画 F(3) 的栈，F(3) 入栈，F(2) 入栈，F(1)
  入栈、出栈，F(0) 入栈、出栈，F(2)
  出栈\ldots\ldots 栈不是一直进，也有出，最深是 \(O(n)\)。
\end{itemize}

老师说本题不需要了解这么深，知道非递归通常效率高就行。

\end{videobox}

\begin{jiexibox}

王道书第 97
页：通常情况下，递归算法在计算机实际执行的过程中包含很多的重复计算，所以效率会低。

\end{jiexibox}

\begin{addbox}

\textbf{老师说“可以当结论”的两个复杂度，补上证明}。用正课 3.3.3 例 2 的
\texttt{Fib}（\texttt{Fib(0)=0}，\texttt{Fib(1)=1}，否则返回
\texttt{Fib(n-1)+Fib(n-2)}）。

\textbf{时间：调用次数是指数级的。} 记 \(C(n)\) 为计算 \texttt{Fib(n)}
的调用次数。和 06 题一样，\(C(0)=C(1)=1\)，\(C(n)=1+C(n-1)+C(n-2)\)。

\emph{断言} \(C(n)=2F_{n+1}-1\)（\(F\)
是斐波那契数）。\emph{归纳}：\(n=0\) 时 \(2F_1-1=1\)，\(n=1\) 时
\(2F_2-1=1\)，成立；若对 \(n-1\)、\(n-2\) 成立，则
\(C(n)=1+(2F_n-1)+(2F_{n-1}-1)=2(F_n+F_{n-1})-1=2F_{n+1}-1\)。

再估计 \(F_{n+1}\) 的大小。因为 \(F\)
单调不减，\(F_{n+1}=F_n+F_{n-1}\le2F_n\)，所以 \(F_{n+1}\le2^n\)；又
\(F_{n+1}=F_n+F_{n-1}\ge2F_{n-1}\)，所以
\(F_{n+1}\ge2^{\lfloor n/2\rfloor}\)。于是

\[2^{\lfloor n/2\rfloor+1}-1\ \le\ C(n)\ \le\ 2^{n+1}-1 .\]

右边说明老师的 \(O(2^n)\) 是对的（它是上界）；左边说明调用次数确实随
\(n\)
指数增长，不只是“不超过”指数。慢的原因就是书上说的重复计算：\texttt{Fib(n-2)}
在 \texttt{Fib(n)} 和 \texttt{Fib(n-1)} 下面各算一次，越往下重复越多。

\textbf{空间：\(O(n)\)。} 任一时刻递归工作栈里的，是从 \texttt{Fib(n)}
一路调用到当前函数的那一串“还没返回”的调用，也就是调用树上从根到当前结点的一条路径（已经返回的调用早已出栈）。路径上每往下一层，参数至少减
1，参数又在 \(0\sim n\) 之间，所以最多 \(n+1\)
层；再细一点，\texttt{Fib(1)} 不再调用别的函数，\texttt{Fib(0)} 只由
\texttt{Fib(2)} 调用，路径不会同时含参数 1 和 0，所以最多 \(n\)
层（\(n\ge1\) 时）；而
\texttt{Fib(n)}→\texttt{Fib(n-1)}→\ldots→\texttt{Fib(1)} 这条路径正好
\(n\) 层，能达到。所以栈的最大深度是 \(n\)，空间复杂度
\(O(n)\)。“一共入栈多少次”是 \(C(n)\)，指数级；“最多同时在栈里多少个”是
\(n\)。空间复杂度看的是后者。

\textbf{贯穿的例子 \texttt{Fib(4)}}：\(C(4)=2F_5-1=2\times5-1=9\)；正课
3.3.3 里我补的 \texttt{Fib(4)} 13 步追踪中，依次调用了
\texttt{Fib(4)}、\texttt{Fib(3)}、\texttt{Fib(2)}、\texttt{Fib(1)}、\texttt{Fib(0)}、\texttt{Fib(1)}、\texttt{Fib(2)}、\texttt{Fib(1)}、\texttt{Fib(0)}，正好
9 次；栈最深时是
\texttt{Fib(4)}、\texttt{Fib(3)}、\texttt{Fib(2)}、\texttt{Fib(1)}
四层（不算 \texttt{main}），等于 \(n=4\)。

\end{addbox}

\begin{recallbox}

正课 3.3.3
知识回顾：递归的缺点是\uln{效率低}，太多层可能栈溢出（\uln{空间复杂度＝递归深度}），可能包含很多重复计算（如斐波那契）；可以自定义栈把递归算法改造成非递归算法。

\end{recallbox}

\begin{questionbox}[15]

【2015 统考真题】已知程序如下：

\begin{Shaded}
\begin{Highlighting}[]
\DataTypeTok{int}\NormalTok{ S}\OperatorTok{(}\DataTypeTok{int}\NormalTok{ n}\OperatorTok{)}
\OperatorTok{\{}   \ControlFlowTok{return} \OperatorTok{(}\NormalTok{n}\OperatorTok{\textless{}=}\DecValTok{0}\OperatorTok{)?}\DecValTok{0}\OperatorTok{:}\NormalTok{S}\OperatorTok{(}\NormalTok{n}\OperatorTok{{-}}\DecValTok{1}\OperatorTok{)+}\NormalTok{n}\OperatorTok{;\}}
\DataTypeTok{void}\NormalTok{ main}\OperatorTok{()}
\OperatorTok{\{}\NormalTok{   cout}\OperatorTok{\textless{}\textless{}}\NormalTok{S}\OperatorTok{(}\DecValTok{1}\OperatorTok{);\}}
\end{Highlighting}
\end{Shaded}

程序运行时使用栈来保存调用过程的信息，自栈底到栈顶保存的信息依次对应的是（　）。

A. \texttt{main()}→\texttt{S(1)}→\texttt{S(0)}　　　B.
\texttt{S(0)}→\texttt{S(1)}→\texttt{main()}\\
C. \texttt{main()}→\texttt{S(0)}→\texttt{S(1)}　　　D.
\texttt{S(1)}→\texttt{S(0)}→\texttt{main()}

\end{questionbox}

\hlred{\textbf{答案：A}}

\begin{videobox}

习题讲解字幕 P32：程序从 \texttt{main} 开始，\texttt{main}
先入栈；\texttt{main} 调用 \texttt{S(1)}，\texttt{S(1)}
入栈；\texttt{S(1)} 中 \(1>0\)，调用 \texttt{S(0)}，\texttt{S(0)}
入栈。栈底到栈顶是 \texttt{main}、\texttt{S(1)}、\texttt{S(0)}，选 A。

\end{videobox}

\begin{jiexibox}

王道书第 99
页：递归调用函数时，在系统栈中保存的函数信息需满足先进后出的特点，依次调用了
\texttt{main()}、\texttt{S(1)}、\texttt{S(0)}，所以栈底到栈顶的信息依次是
\texttt{main()}、\texttt{S(1)}、\texttt{S(0)}。

\textbf{注意}：在递归中，系统为每一层的返回点、局部变量、传入实参等开辟了递归工作栈来存储。

\end{jiexibox}

\begin{addbox}

\textbf{栈的变化过程}（左边是栈底）：\texttt{main} →
\texttt{main,\ S(1)} →
\texttt{main,\ S(1),\ S(0)}（最深，本题问的就是此时）→ \texttt{S(0)}
返回 0，出栈：\texttt{main,\ S(1)} → \texttt{S(1)} 返回
\(0+1=1\)，出栈：\texttt{main} → 输出 1。

栈底一定是最先被调用、最后结束的 \texttt{main}，所以 B、D 一眼可排除；C
把调用次序弄反了，\texttt{S(0)} 是 \texttt{S(1)} 调用的，只能在它上面。

\end{addbox}

\begin{recallbox}

正课 3.3.3
第一节：函数调用的特点是\uln{最后被调用的函数最先执行结束}；栈底在
\texttt{main}
下面还有编译器加的一些信息，题目里不考虑。书上“返回点、局部变量、传入实参”＝正课函数调用栈的三样东西。

\end{recallbox}

\hypertarget{ux56dbux5173ux4e8eux6808ux7684ux5224ux65ad1117}{%
\section{四、关于栈的判断（11、17）}\label{ux56dbux5173ux4e8eux6808ux7684ux5224ux65ad1117}}

\tocback

\Needspace{8\baselineskip}\vspace{22pt}

\cueblock{\cuenote}{\relax{判断“一定”“必须”“只要\ldots\ldots 就”时，用什么证明它错？}}

\begin{questionbox}[11]

下列说法中，正确的是（　）。

A. 消除递归不一定需要使用栈\\
B.
对同一输入序列进行两组不同的合法入栈和出栈组合操作，所得的输出序列也一定相同\\
C. 通常使用队列来处理函数或过程调用\\
D. 队列和栈都是运算受限的线性表，只允许在表的两端进行运算

\end{questionbox}

\hlred{\textbf{答案：A}}

\begin{videobox}

习题讲解字幕 P32：A 对，斐波那契数列用循环就能算，不用栈。B 错，入栈
1、2、3，可以得到 2、1、3，也可以得到 3、2、1。C 错，函数调用用栈。D
错，栈只在一端（栈顶）操作。

\end{videobox}

\begin{jiexibox}

王道书第 97--98
页：使用栈可以模拟递归的过程，以此来消除递归，但对于单向递归和尾递归而言，可以用迭代的方式来消除递归，选项
A 正确。不同的入栈和出栈组合操作，会产生许多不同的输出序列，B
错误。通常使用栈来处理函数或过程调用，选项 C
错误。队列和栈都是操作受限的线性表，但只有队列允许在表的两端进行运算，而栈只允许在栈顶方向进行操作，选项
D 错误。

\end{jiexibox}

\begin{addbox}

\textbf{每个选项的证明各需要什么}：

\begin{itemize}
\tightlist
\item
  A 说“不一定需要”，否定的是“消除递归一定要用栈”。要证明
  A，\uln{找到一个不用栈就能消除的递归}就够了：08
  题老师写的斐波那契循环，只用
  \texttt{pre1}、\texttt{pre2}、\texttt{current} 三个变量，没有栈；
\item
  B 说“一定相同”，是全称说法，\uln{一个反例就证明它错}：入栈序列 1, 2,
  3，操作“入 入 出 入 出 出”输出 2, 1, 3，“入 入 入 出 出 出”输出 3, 2,
  1；
\item
  C 与正课 3.3.3 矛盾（函数调用用栈）；
\item
  D 中“栈允许在两端运算”与栈的定义矛盾。
\end{itemize}

\end{addbox}

\begin{defbox}[：名词统一]

书上“单向递归”“尾递归”正课没有讲。尾递归指递归调用是函数的最后一步（调用返回后不再做别的事），比如正课
3.3.3
的阶乘可以改写成这种形式；这种递归每层都不用保留信息，用一个循环就能代替。

\end{defbox}

\begin{questionbox}[17]

【2017 统考真题】下列关于栈的叙述中，错误的是（　）。

Ⅰ. 采用非递归方式重写递归程序时必须使用栈\\
Ⅱ. 函数调用时，系统要用栈保存必要的信息\\
Ⅲ. 只要确定了入栈次序，即可确定出栈次序\\
Ⅳ. 栈是一种受限的线性表，允许在其两端进行操作

A. 仅 Ⅰ　　B. 仅 Ⅰ、Ⅱ、Ⅲ　　C. 仅 Ⅰ、Ⅲ、Ⅳ　　D. 仅 Ⅱ、Ⅲ、Ⅳ

\end{questionbox}

\hlred{\textbf{答案：C}}

\begin{videobox}

习题讲解字幕 P33：Ⅰ
错，还是斐波那契的例子，用循环就能打破递归形式，不一定要用栈。Ⅱ
对，函数调用时系统把相关信息压进函数栈。Ⅲ 错，入栈次序 1, 2,
3，全部进栈再全部出栈得 3, 2, 1；进一个出一个得 1, 2,
3。老师特别提醒：\uln{如果把 Ⅲ 里的栈换成队列，那就对了}，确定了入队次序就确定了出队次序，下次改成队列要注意。Ⅳ
错，栈只允许在一端操作，进栈出栈都靠栈顶指针。错误的是 Ⅰ、Ⅲ、Ⅳ，选 C。

\end{videobox}

\begin{jiexibox}

王道书第 99 页：说法 Ⅰ
的反例：计算斐波那契数列迭代实现只需要一个循环即可实现。说法 Ⅲ
的反例：入栈序列为 1, 2，进行 Push, Push, Pop, Pop 操作，出栈次序为
2、1；进行 Push, Pop, Push, Pop 操作，出栈次序为 1, 2。说法
Ⅳ，栈是一种受限的线性表，只允许在一端进行操作。说法 Ⅱ 正确。

\end{jiexibox}

\begin{addbox}

\textbf{老师说“换成队列就对了”，证明一下}：队列只在队尾插入、队头删除，任何两个元素
\(x\) 先于 \(y\) 入队，\(y\) 排在 \(x\) 后面；\(y\)
要出队必须先到队头，而 \(x\) 在它前面，所以 \(x\)
一定先出队。对任意两个元素都成立，所以出队次序和入队次序完全相同，只有一种。栈没有这个性质：\(x\)
先入栈后，\(y\) 入栈时 \(x\) 可能已经出栈了（\(x\)
先出），也可能还在（\(y\) 先出），两种都合法，这就是 Ⅲ 的反例。

这和 3.2 课后习题 04 的结论一样：一个入队序列只有一种出队序列。

\end{addbox}

\begin{recallbox}

3.1 课后习题的出栈序列题：同一入栈序列有多种出栈序列，个数是卡特兰数。11
的 B、17 的 Ⅲ 都是这个事实。

\end{recallbox}

\hypertarget{ux4e94ux62ecux53f7ux5339ux914dux548cux706bux8f66ux8f68ux90532016}{%
\section{五、括号匹配和火车轨道（20、16）}\label{ux4e94ux62ecux53f7ux5339ux914dux548cux706bux8f66ux8f68ux90532016}}

\tocback

\Needspace{8\baselineskip}\vspace{22pt}

\cueblock{\cuenote}{\relax{括号匹配时栈最深多少，和括号的什么有关？火车轨道为什么至少 4 条，凭什么
4 条就够？}}

\begin{questionbox}[20]

【2025 统考真题】已知算法 A 用于检查字符串中各类括号是否匹配，A
执行过程中使用初始为空的栈保存遇到的括号。若栈的容量是
3，则下列选项中，A 不能处理的是（　）。

A. \texttt{(a+{[}b+(c+d)/e{]}+f)+g-h}　　　B.
\texttt{{[}a*((b+c)/(d-e)+f/g)-h{]}}\\
C. \texttt{{[}a*(b-(c-d)*e/(f+g))-h{]}}　　　D.
\texttt{{[}a-(b+{[}c*(d+e)-f{]}+g+h){]}}

\end{questionbox}

\hlred{\textbf{答案：D}}

\begin{videobox}

习题讲解字幕
P33（一休老师）：这道题比较仁慈，只需要处理括号，括号以外的东西可以视而不见。先把每个选项只保留括号，再模拟
D：\texttt{{[}} 入栈，\texttt{(} 入栈，\texttt{{[}} 入栈，\texttt{(}
还要入栈，这时栈里要放 4 个，超过容量 3，处理不了。

更简单的做法：遇到左括号入栈、右括号出栈，所以设一个计数器记栈里的元素个数，\uln{左括号加一，右括号减一}，看最大值。A、B、C
最大是 3，D 是 4，选 D。

老师顺便说：中缀表达式的考法主要有两种，中缀转后缀（模拟，更重要）和用有向图描述表达式。

\end{videobox}

\begin{jiexibox}

王道书第 99--100
页：该算法利用栈来检查括号匹配：遇到左括号时入栈，遇到右括号时出栈并与之匹配。由于栈的容量为
3，因此最多只能容纳 3 个未匹配的左括号，即最多支持 3
层嵌套的括号结构。逐项分析各选项的括号嵌套深度：对于选项 A，括号序列为
\texttt{({[}(){]})}，最大嵌套深度为 3，依次入栈
\texttt{({[}(}，未超过栈容量。对于选项 B，括号序列为
\texttt{{[}(()){]}}，最大深度为 3，依次入栈
\texttt{{[}((}，栈容量足够。对于选项 C，括号序列为
\texttt{{[}(()){]}}，同样最大深度为 3，栈容量足够。对于选项
D，括号序列为 \texttt{{[}({[}(){]}){]}}，最大嵌套深度达到 4。扫描至
\texttt{*} 后的 \texttt{(} 时，栈中已存有 \texttt{{[}({[}}
三个左括号，此时，需要将第四个左括号 \texttt{(} 入栈，但栈容量仅为
3，无法继续入栈，因此该表达式无法被正确处理。

\end{jiexibox}

\begin{addbox}

\textbf{计数器为什么就等于栈里的个数（证明）}：正课 3.3.1
的算法，栈里只放左括号；遇到左括号压入一个，遇到右括号弹出一个（匹配成功时）。所以在括号一直匹配成功的前提下，扫描完某个前缀时，

\[\text{栈中元素个数} = (\text{这个前缀里的左括号数}) - (\text{右括号数}),\]

也就是此处的括号嵌套层数。容量为 3
的栈能处理整个串，\uln{当且仅当所有前缀的这个差都不超过 3}。

\textbf{逐个选项写出计数}（只看括号）：

{\rowcolors{2}{wdblue!5}{white}\arrayrulecolor{wdblue!35}
\begin{xltabular}{\linewidth}{>{\hsize=0.400\hsize\raggedright\arraybackslash}X>{\hsize=0.880\hsize\raggedright\arraybackslash}X>{\hsize=1.920\hsize\raggedright\arraybackslash}X>{\hsize=0.800\hsize\raggedright\arraybackslash}X}
\rowcolor{wdblue}\textcolor{white}{\bfseries 选项} & \textcolor{white}{\bfseries 括号序列} & \textcolor{white}{\bfseries 每个括号之后的计数} & \textcolor{white}{\bfseries 最大} \\
\endfirsthead
A & \texttt{(\ {[}\ (\ )\ {]}\ )} & 1, 2, 3, 2, 1, 0 & 3 \\ \hline
B & \texttt{{[}\ (\ (\ )\ (\ )\ )\ {]}} & 1, 2, 3, 2, 3, 2, 1, 0 & 3 \\ \hline
C & \texttt{{[}\ (\ (\ )\ (\ )\ )\ {]}} & 1, 2, 3, 2, 3, 2, 1, 0 & 3 \\ \hline
D & \texttt{{[}\ (\ {[}\ (\ )\ {]}\ )\ {]}} & 1, 2, 3, \hlred{4}, 3, 2, 1, 0 & \hlred{4} \\ \hline
\end{xltabular}\par\vspace{4pt}}

书上把 B、C 写成 \texttt{{[}(()){]}}，少写了一对：B 的
\texttt{((b+c)/(d-e)+f/g)} 里有 \texttt{(b+c)}、\texttt{(d-e)} 两对，C
的 \texttt{(b-(c-d)*e/(f+g))} 里有 \texttt{(c-d)}、\texttt{(f+g)}
两对，都是 \texttt{{[}(()()){]}}。最大深度仍是 3，不影响答案。

\end{addbox}

\begin{recallbox}

正课
3.3.1：\uln{最后出现的左括号最先被匹配}，所以用栈；顺序栈容量固定，括号串很长时可能栈溢出。

\end{recallbox}

\begin{questionbox}[16]

【2016 统考真题】设有如下图所示的火车车轨，入口和出口之间有 n
条轨道，列车的行进方向均为从左至右，列车可驶入任意一条轨道。现有编号为
1～9 的 9 列列车，驶入的次序依次是 8, 4, 2, 5, 3, 9, 1, 6,
7。若期望驶出的次序依次为 1～9，则 n 至少是（　）。

A. 2　　　B. 3　　　C. 4　　　D. 5

\end{questionbox}

\slidefig{assets/3.3ex-16-book.png}{0.75}{王道书第 96 页：16 题的火车车轨}

\hlred{\textbf{答案：C}}

\begin{videobox}

习题讲解字幕
P33：题目特意标出了入口和出口，一边进一边出，每条轨道就是一个队列，问的是至少要几个队列。逐个安排：8
先随便进一条；4 要比 8 先出，不能排在 8 后面，只能另开一条；2 要比 4、8
都先出，再开一条；5 可以排在 4 后面；3 排在 2 后面；9 排在 8 后面；1
要第一个出，只能单独再开一条；6、7 要在 4、5 之后出，也在 2、3
之后，放在哪条都行。共 4 条，选 C。

老师说：“这种题目都是要求我们手动去模拟，并没有什么严格上的一个理论证明的东西。”

\end{videobox}

\begin{jiexibox}

王道书第 99 页：根据题意：入队顺序为 8, 4, 2, 5, 3, 9, 1, 6,
7，出队顺序为 1～9。入口和出口之间有多个队列（n
条轨道），且每个队列（轨道）可容纳多个元素（多列列车），为便于区分，队列用字母编号。分析如下：显然先入队的元素必须小于后入队的元素（否则，若
8 和 4 入同一个队列，8 在 4 前面，则出队时也只能 8 在 4 前面），这样 8
入队列 A，4 入队列 B，2 入队列 C，5 入队列
B（按照前述原则“大的元素在小的元素后面”也可将 5 入队列
C，但这时剩下的元素 3 就必须放入一个新的队列中，无法确保“至少”），3
入队列 C，9 入队列 A，这时共占了 3 个队列，后面还有元素
1，直接再用一个新的队列 D，1 从队列 D 出队后，剩下的元素 6 和 7 或入队列
B，或入队列 C。综上，共占用了 4
个队列。当然还有其他的入队、出队情况，请读者自己推演，但要确保满足：①
队列中后面的元素大于前面的元素；②
确保占用最少（满足题意中“至少”）的队列。

\end{jiexibox}

\begin{addbox}

\textbf{这道题其实有严格证明}：老师和书都只给出了一种 4
条轨道的安排。但“至少是
4”包含两句话，要分别证明：\uln{3 条不够}（下界），\uln{4 条够}（构造）。

\textbf{引理（同一条轨道上的车必须按编号从小到大驶入）}：轨道是队列，同一条轨道上
\(x\) 先于 \(y\) 驶入，\(x\) 就先于 \(y\) 驶出（17
题我补的证明）。出口要求编号小的先出，所以同一条轨道上先驶入的车编号必须更小。

\textbf{下界，3 条不够}：在驶入次序 8, 4, 2, 5, 3, 9, 1, 6, 7
里，8、4、2、1
依次出现、编号依次变小。由引理，其中任意两列都不能在同一条轨道上（先进来的编号大，违反引理）。4
列车两两不同轨，至少要 4 条轨道。

\textbf{构造，4 条够}：按书上的安排（6、7 都放 B）：

\end{addbox}

下面是驶入、驶出全过程的分解图（我画的）。每一小格是一步：每条轨道右端靠出口，是最前面的车；\uln{绿色是这一步驶入的车}；右边是已驶出的车，\hlred{红色是这一步驶出的}。

\slidefig{assets/3.3ex-16-train-a.png}{1.0}{16 题：9 列车依次驶入（第 1--8 步）}

\slidefig{assets/3.3ex-16-train-b.png}{1.0}{16 题：最后一列车驶入，然后依次驶出 1～9（第 9--14 步）}

\begin{addbox}

\textbf{速查}（接上）：上面的过程合成一张表。

{\rowcolors{2}{wdblue!5}{white}\arrayrulecolor{wdblue!35}
\begin{xltabular}{\linewidth}{>{\hsize=0.600\hsize\raggedright\arraybackslash}X>{\hsize=1.200\hsize\raggedright\arraybackslash}X>{\hsize=1.200\hsize\raggedright\arraybackslash}X}
\rowcolor{wdblue}\textcolor{white}{\bfseries 轨道} & \textcolor{white}{\bfseries 车（按驶入先后）} & \textcolor{white}{\bfseries 是否从小到大} \\
\endfirsthead
A & 8, 9 & 是 \\ \hline
B & 4, 5, 6, 7 & 是 \\ \hline
C & 2, 3 & 是 \\ \hline
D & 1 & 是 \\ \hline
\end{xltabular}\par\vspace{4pt}}

驶出：让所有车先驶入，然后每次从 4
条轨道的最前面（队头）里选编号最小的驶出。因为每条轨道都是从小到大排的，每条轨道的队头就是这条轨道上最小的车，所以所有队头中最小的，就是全体剩下的车中最小的。于是第
1 次驶出的是 1，第 2 次是 2\ldots\ldots 依次正好驶出 1～9。

\textbf{贯穿这个例子看驶出过程}（每行是驶出前各轨道的队头）：

{\rowcolors{2}{wdblue!5}{white}\arrayrulecolor{wdblue!35}
\begin{xltabular}{\linewidth}{>{\hsize=0.480\hsize\raggedright\arraybackslash}X>{\hsize=1.260\hsize\raggedright\arraybackslash}X>{\hsize=1.260\hsize\raggedright\arraybackslash}X}
\rowcolor{wdblue}\textcolor{white}{\bfseries 第几次} & \textcolor{white}{\bfseries 队头 A / B / C / D} & \textcolor{white}{\bfseries 驶出} \\
\endfirsthead
1 & 8 / 4 / 2 / 1 & 1（D 空） \\ \hline
2 & 8 / 4 / 2 / -- & 2 \\ \hline
3 & 8 / 4 / 3 / -- & 3（C 空） \\ \hline
4--7 & 8 / 4,5,6,7 依次 / -- / -- & 4, 5, 6, 7（B 空） \\ \hline
8--9 & 8, 9 / -- / -- / -- & 8, 9 \\ \hline
\end{xltabular}\par\vspace{4pt}}

两部分合起来，\(n\) 的最小值就是 4。

\textbf{推广}：同样的道理说明，所需轨道数至少是驶入序列中“编号依次变小的最长子序列”的长度。本题这个长度是
4（8, 4, 2, 1），所以下界是 4；而上面又构造出了 4
条的方案，两者相等。考场上找出这条最长下降子序列，就知道答案的下界；再像老师那样贪心地安排一遍，安排出来的条数等于下界时，答案就定了。

\end{addbox}

\begin{recallbox}

3.2 课后习题
04：队列先进先出，\uln{一个入队序列只有一种出队序列}；本题把这个性质用在每一条轨道上。正课
3.3.4+3.3.5
书框中的考点追踪“多队列出队/入队操作的应用（2016）”就是这道题。

\end{recallbox}

\newpage

\hypertarget{ux516dux603bux7ed3}{%
\section{六、总结}\label{ux516dux603bux7ed3}}

\tocback

\begin{summarybox}

\textbf{用栈还是用队列}（01、03、09、10、12）：后来的先处理用栈（括号匹配、表达式求值、递归和函数调用、局部变量、树的递归遍历、DFS）；先来的先处理用队列（缓冲区、FIFO
页面替换、层次遍历、BFS、先来先服务）。

\textbf{表达式}（02、19、04、13、14、18）：

\begin{itemize}
\tightlist
\item
  求后缀：找最后执行的运算符（括号外、优先级最低、同级取最右）拆成 \(L\)
  op \(R\)，后缀＝\(L\) 的后缀 \(R\) 的后缀
  op；它等于表达式树的后序遍历，“\(F_1\)
  最长”就是同级从左往右算（02、19）。后缀唯一，所以排除法可靠；
\item
  运算数栈深度：转成后缀，从左往右数“操作数个数 −
  运算符个数”的最大值；或者在每个操作数处数“未运算的运算符（不算括号）＋1”（04）；
\item
  中缀转后缀栈中的情况：完整模拟一遍，括号也占位置（13）；扫描到某个操作数时，栈中运算符＝中缀式中在它之前、后缀式中在它之后的运算符，再补上未消解的
  \texttt{(}（14）；
\item
  先弹出的是右操作数，b op a（18）。
\end{itemize}

\textbf{递归}（05--08、15）：画递归调用树。调用次数＝结点数，可以用递推式
\(C(n)=1+\ldots\)
算，不要漏了最外层那次（06）；执行次序＝先序遍历（07）；栈里是根到当前结点的路径，最大深度＝树高（15）。斐波那契递归调用
\(2F_{n+1}-1\) 次，指数级；空间
\(O(n)\)，看的是“最多同时在栈里多少”，不是“一共入栈多少次”（08）。

\textbf{判断题}（11、17）：否定“一定”“必须”“只要\ldots\ldots 就”，一个反例就够（斐波那契循环、入栈
1,2 的两种出栈次序）。“入队次序确定出队次序”对队列成立，对栈不成立。

\textbf{括号匹配栈深}（20）：栈中个数＝前缀里左括号数 −
右括号数＝嵌套层数；容量够当且仅当最大嵌套层数不超过容量。

\textbf{多队列}（16）：同一队列里必须从小到大。驶入序列中最长的下降子序列每个元素都要单独一条，给出下界；再构造一个方案达到它。

\textbf{自测}：

\begin{enumerate}
\def\labelenumi{\arabic{enumi}.}
\tightlist
\item
  写出 \texttt{a-(b+c*d)/e}
  的后缀表达式，并说出每一步拆分时最后执行的是哪个运算符。
\item
  用两个栈计算 \texttt{a-(b+c*d)/e} 时，运算数栈最深几个？用“操作数 −
  运算符”的方法算。
\item
  13 题中如果把栈容量限制为 4，转换会在扫描到哪个字符时溢出？
\item
  06 题的 F 改成
  \texttt{if(n\textless{}=3)\ return\ 1;\ else\ return\ F(n-1)+F(n-3)+1;}，计算
  F(6) 调用几次？
\item
  为什么 08 题递归求斐波那契的空间复杂度是 \(O(n)\) 而不是 \(O(2^n)\)？
\item
  16 题如果驶入次序是 9, 8, 7, 6, 5, 4, 3, 2,
  1，至少要几条轨道？为什么？
\end{enumerate}

\end{summarybox}
\end{document}
